% Changing the ATE application - the developer's manual (packet 34c).
%
% Build: python docs/manuals/build.py developer
% (latexmk -lualatex inside texlive/texlive:latest; see docs/building-from-source.md.)
%
% Every fact here is checked against the code. tests/test\_manuals.py derives the module list, the
% Qt signals, the schema tables and the Device drivers from the code and fails if this manual or
% docs/manuals/llm/developer.md stops naming them; every command in a listing is either run by the
% suite or listed with the reason it is not; the two worked examples (a Device driver, a schema
% version) are executed against a temporary copy of the source. Where a document disagrees with
% the code, the code is followed and the disagreement is in the last chapter.

\documentclass[a4paper,11pt]{report}

\usepackage{manual}  % docs/manuals/common/manual.sty, on TEXINPUTS by the build script
\usepackage{longtable}  % the module table in chapter 2 is longer than one page

% Long code names are what makes a line here overflow: the build gate fails on any box more than
% 5pt too wide, so line breaking gets extra room, and every path or dotted name in the text is
% written with \allowbreak{} where it can be cut.
\setlength{\emergencystretch}{2em}
\sloppy
% The developer manual quotes shell commands as well as Lua, so it needs the shell listing style
% in addition to the Lua one the shared preamble defines.
\lstdefinelanguage{ATEShell}{
  morekeywords={ate,python,pytest,latexmk,docker,git,apt,pip,venv,source,export,lua,cd,mkdir,cp,ls,latex,find,grep},
  sensitive=true,
  morecomment=[l]{\#},
  basicstyle=\ttfamily\footnotesize,
  keywordstyle=\bfseries\color{codekey},
  commentstyle=\itshape\color{codecomment},
  backgroundcolor=\color{codebg},
  frame=single,
  rulecolor=\color{ategray},
  breaklines=true,
  breakatwhitespace=true,
  captionpos=b,
  numbers=none,
  columns=fullflexible,
  keepspaces=true,
  showstringspaces=false
}

% Python listings: the same shape as the Lua one, with Python's keywords.
\lstdefinelanguage{ATEPython}{
  keywords={def,class,return,if,then,else,elif,end,for,while,import,from,as,with,not,and,or,None,True,False,lambda,try,except,finally,raise,pass,assert,in,is,global,yield,self},
  sensitive=true,
  morecomment=[l]{\#},
  basicstyle=\ttfamily\footnotesize,
  keywordstyle=\bfseries\color{codekey},
  commentstyle=\itshape\color{codecomment},
  stringstyle=\color{ategreen},
  backgroundcolor=\color{codebg},
  frame=single,
  rulecolor=\color{ategray},
  breaklines=true,
  breakatwhitespace=true,
  captionpos=b,
  numbers=none,
  columns=fullflexible,
  keepspaces=true,
  showstringspaces=false
}

\title{%
  \textbf{Changing the ATE Application}\\[6pt]
  \large The developer's manual\\[10pt]
  \normalsize ATE v1.0 \quad\textbullet\quad every name here is checked against the code by the test suite}
\author{ATE}
\date{November 2026}

\begin{document}
\maketitle
\tableofcontents

\chapter{What this manual is, and what it has proven}\label{ch:intro}

This manual is for the person who has to change the ATE application: a Python developer who may not
have written a Test Schedule, and who may never have threaded a Qt application before. The rules here
are not style rules. Each one exists because a run broke in a specific way, and the fix is a rule
about where code may go and what it may call. A rule you can recite but not explain is a rule the next
change breaks.

Two files make this manual, and they say the same things:

\begin{itemize}
  \item this PDF, built from \texttt{docs/manuals/developer/} by \texttt{python docs/manuals/build.py developer};
  \item \texttt{docs/manuals/llm/developer.md}, written for an agent rather than converted from this
    text: it is shorter, it is in the order an agent works, and the checks below apply to it too.
\end{itemize}

\section{What is checked, and how}

\texttt{tests/test_manuals.py} holds the manuals to the code. It does not read them for style; it
derives lists of names from the code and fails when a manual stops naming one:

\begin{itemize}
  \item every module under \texttt{src/ate/}, derived with \texttt{pkgutil} from the package itself;
  \item every Qt signal declared on \texttt{RunEngine} or \texttt{HeadlessRunSession}, read off the
    classes;
  \item every table in \texttt{ate.storage.schema}, read from the SQLAlchemy \texttt{metadata};
  \item every Device driver class and the Device Type name it serves, taken from the classes that
    subclass \texttt{ate.devices.registry.Driver}.
\end{itemize}

The same test extracts every \texttt{src/...} path and every dotted name in backticks from the two
developer manuals and looks them up: a file that does not exist, or a module that does not have the
attribute named, fails the suite. Commands in a shell listing are each either run by the test suite in
a temporary Station or listed with the reason they cannot be run here. The two worked examples,
adding a Device type (chapter~\ref{ch:device-type}) and adding a schema version
(chapter~\ref{ch:schema-version}), are executed against a temporary copy of the source: the steps the
manual gives are the steps the test performs.

\section{How to run things}

Nothing here needs a display. Every Qt test runs with the offscreen platform:

\begin{lstlisting}[language=ATEShell]
QT_QPA_PLATFORM=offscreen python -m pytest -q -p no:cacheprovider
\end{lstlisting}

That is the command CI runs, and the one this manual was written against. Two gates are off by
default because they take minutes or need Docker:

\begin{lstlisting}[language=ATEShell]
ATE_BUILD_TESTS=1 pytest -m build      # builds the real PyInstaller bundle
ATE_DOCKER_TESTS=1 pytest -m build     # runs the debian:stable container smokes
\end{lstlisting}

Both are described in chapter~\ref{ch:tests}, and both were run for this manual.

\section{What is proven here, and what is not}

Proven on this machine, by the tests named in the chapters that state them:

\begin{itemize}
  \item the whole suite on Linux with the offscreen platform, repeated until it was green;
  \item the threading behaviour of every signal the engine and the session emit, and the thread each
    slot runs on (\texttt{tests/test_signal_threads_32a.py});
  \item that an abort ends a Device call on every Device Type, and how long it takes
    (\texttt{tests/test_device_abort_32b.py});
  \item that a new Device driver and a new schema version work when added the way this manual says.
\end{itemize}

Not proven, and this manual says so wherever it matters:

\begin{itemize}
  \item \textbf{real hardware.} Every Device test here talks to a loopback peer in
    \texttt{tests/_fixtures.py} or a pseudo-terminal. A GPIB instrument, a real serial port and a
    real VISA stack have never been touched by the suite.
  \item \textbf{real Windows sessions.} Windows behaviour is asserted by the same tests, but only
    GitHub Actions has ever run them on Windows, and only as CI runs.
  \item \textbf{the cause of the 32b stall.} One Windows CI run stalled a Test Run for tens of
    seconds. The fixes are proven; the cause is a belief about the disk, not a measurement.
    Chapter~\ref{ch:failed-run} says what a future failure would show.
\end{itemize}

\cando{tell which claim in this manual rests on a test you can run, and which rests on a machine that
is not here.}

\chapter{The shape of the code}\label{ch:shape}

ATE is one process, one \texttt{QApplication}, and one worker thread per Test Run. The source under
\texttt{src/ate/} is small enough to read in an afternoon, and the structure is the thing that makes it
safe to change: a change in one layer cannot reach another layer except through a signal or a store
call.

\section{Four layers, one direction}

The dependency arrows go one way:

\begin{itemize}
  \item \texttt{ate.ui} knows \texttt{ate.run_session}. A layout connects to the session's signals and
    calls its methods. It never reaches into the engine.
  \item \texttt{ate.run_session} knows \texttt{ate.engine}. The session owns a \texttt{RunEngine},
    forwards its signals, and writes nothing to the database itself.
  \item \texttt{ate.engine} knows \texttt{ate.storage} and \texttt{ate.devices}. The worker writes rows
    and opens Devices. It knows nothing about screens.
  \item \texttt{ate.storage} is the only place SQL is written. \texttt{ate.catalog}, the configuration
    reader and the report writer are reached from above, not from the engine.
\end{itemize}

Nothing points the other way. \texttt{ate/__init__.py} deliberately imports neither the engine nor the
UI, so \texttt{import ate} stays cheap and widget-free: it exposes only \texttt{__version__} and
\texttt{version}.

\section{Every module}

This table is the coverage list: \texttt{tests/test_manuals.py} walks the package with
\texttt{pkgutil.walk_packages} and fails if a module in it is not named here or in
\texttt{docs/manuals/llm/developer.md}.

\begin{center}
\small
\begin{tabular}{p{4.6cm} p{7.4cm}}
\toprule
\textbf{File} & \textbf{What it holds} \\
\midrule
\texttt{src/ate/\_\_init\_\_.py} & the package's version only; imports no engine, no widgets \\
\texttt{src/ate/\_version.py} & the version from the git tag, computed once so build time and run time agree \\
\texttt{src/ate/\_\_main\_\_.py} & \texttt{python -m ate}, the entry the PDF helper subprocess uses \\
\texttt{src/ate/cli.py} & the \texttt{ate} command: parsing, subcommands, \texttt{ate setup}, \texttt{ate db}, the hidden \texttt{ate \_pdf} and \texttt{ate \_selftest} \\
\texttt{src/ate/app.py} & the application shell: reads config once, opens the chosen layout, installs the exception hook \\
\texttt{src/ate/config.py} & the config file: read, validate, fill defaults, resolve \texttt{<data>} \\
\texttt{src/ate/permissions.py} & the group, ownership and ACL treatment \texttt{ate setup} applies on Linux \\
\texttt{src/ate/setup.py} & prepare a Station: directories, default config, database, permissions \\
\texttt{src/ate/log.py} & the per-launch application log, \texttt{AppLog} \\
\texttt{src/ate/errors.py} & \texttt{ATEError} and its subclasses; the one base every ATE failure has \\
\texttt{src/ate/exception_log.py} & \texttt{UnhandledExceptionLog}: what a raising slot writes to the log \\
\texttt{src/ate/catalog.py} & \texttt{.atedef} reading and validation, Media Release chains, schedule and module maps \\
\texttt{src/ate/engine/\_\_init\_\_.py} & the engine package; the rule is that it imports no widgets \\
\texttt{src/ate/engine/run_engine.py} & the \texttt{QThread} worker: steps the schedule, records the run, ends it \\
\texttt{src/ate/engine/lua_host.py} & the Lua runtime, the sandbox, the \texttt{ATE} table, the schedule builders \\
\texttt{src/ate/engine/operator_input.py} & \texttt{OperatorInput}: how \texttt{ATE.input(n)} waits for an answer \\
\texttt{src/ate/engine/step_inputs.py} & parsing and validating a Step Input declaration \\
\texttt{src/ate/lua/\_\_init\_\_.py} & the modules bundled with ATE, reachable as \texttt{require("ate.devices")} \\
\texttt{src/ate/lua/devices.lua} & the Lua module a Test Schedule can \texttt{require} \\
\texttt{src/ate/run_session/\_\_init\_\_.py} & the run session package; also widget-free \\
\texttt{src/ate/run_session/headless.py} & \texttt{HeadlessRunSession}: the object a layout binds to, and what a headless run uses \\
\texttt{src/ate/storage/\_\_init\_\_.py} & the storage package: the only place ATE talks SQL \\
\texttt{src/ate/storage/schema.py} & every table in the Test Run record, as SQLAlchemy Core tables \\
\texttt{src/ate/storage/store.py} & \texttt{RunStore}: writes, queries, batching, the schema check \\
\texttt{src/ate/storage/\_uuid.py} & UUIDv7 run keys, without \texttt{uuid.uuid7} \\
\texttt{src/ate/devices/\_\_init\_\_.py} & the Device drivers package \\
\texttt{src/ate/devices/registry.py} & \texttt{Driver}, \texttt{validate\_params}, \texttt{DeviceRegistry} \\
\texttt{src/ate/devices/stream.py} & the background reader shared by the buffered Device Types \\
\texttt{src/ate/devices/tcp.py} & the TCP driver \\
\texttt{src/ate/devices/udp.py} & the UDP driver \\
\texttt{src/ate/devices/serial.py} & the SERIAL driver, through pyserial \\
\texttt{src/ate/devices/visa.py} & the VISA driver, through pyvisa \\
\texttt{src/ate/devices/abort.py} & how a Device call gives way to an abort \\
\texttt{src/ate/devices/tracing.py} & one \texttt{DEBUG} line per Device operation, when tracing is on \\
\texttt{src/ate/report/\_\_init\_\_.py} & the report package \\
\texttt{src/ate/report/html.py} & the self-contained HTML Test Report \\
\texttt{src/ate/report/pdf.py} & HTML to PDF in a separate process, and the sandbox decision \\
\texttt{src/ate/report/pdf_job.py} & the \texttt{QProcess} that runs the helper, owned by a session \\
\texttt{src/ate/report/log_export.py} & the Run log as CSV, JSON Lines or text \\
\texttt{src/ate/history/\_\_init\_\_.py} & the history package \\
\texttt{src/ate/history/session.py} & \texttt{HistorySession}: queries over past Test Runs, no widgets \\
\texttt{src/ate/ui/\_\_init\_\_.py} & the widgets package; the only place QtWidgets is imported \\
\texttt{src/ate/ui/run_widgets.py} & the widgets both Run screens share \\
\texttt{src/ate/ui/staged_run_screen.py} & the \texttt{staged} layout \\
\texttt{src/ate/ui/three_pane_run_screen.py} & the \texttt{threePane} layout \\
\texttt{src/ate/ui/history_screen.py} & the History screen \\
\texttt{src/ate/ui/export_log_dialog.py} & the Run log export dialog \\
\texttt{src/ate/dialogs.py} & \texttt{dismiss_if_headless}: every modal that must answer itself \\
\bottomrule
\end{tabular}
\end{center}

Outside \texttt{src/} there is \texttt{packaging/} (PyInstaller, Inno Setup, the \texttt{.deb}, the wheel
and the licence inventory, chapter~\ref{ch:packaging}), \texttt{examples/} (a Media Release and the
harness that runs it, used by the test suite), and \texttt{docs/} (the ADRs, this manual and the agent
manuals).

\section{What a Test Run is made of, in the code}

Four objects move one Test Run through the application, and knowing which one you are in is half of
debugging it:

\begin{itemize}
  \item the \textbf{run session} (\texttt{HeadlessRunSession}) holds the Operator's choices: which
    schedule, which Operator, whether Stop on Fail is on. It is a \texttt{QObject}, so widgets can bind
    to it.
  \item the \textbf{run engine} (\texttt{RunEngine}, a \texttt{QThread}) is the worker that walks the
    expanded step tree and writes the record.
  \item the \textbf{Lua host} (\texttt{LuaHost}) is the interpreter the schedule runs in, with the
    sandbox and the \texttt{ATE} table.
  \item the \textbf{store} (\texttt{RunStore}) is the only thing that touches the database.
\end{itemize}

A layout is a fifth object, above the session: \texttt{StagedRunScreen} or \texttt{ThreePaneRunScreen},
chosen by \texttt{ui.layout} and switched in \texttt{App.set\_layout}.

\cando{say which object owns a piece of state, and which layer a new line of code belongs in.}

\chapter{The rule the code protects: no widgets in the engine}\label{ch:rule}

The spec's central architectural rule is that the engine and the run session never import widgets, so
a Test Run can be driven, tested and run with no display at all. That rule is enforced by a test, not
by review: \texttt{tests/test_architecture_rules.py} parses every \texttt{.py} under
\texttt{src/ate/engine} and \texttt{src/ate/run_session} with \texttt{ast} and fails if any of them
imports \texttt{QtWidgets}, \texttt{uic}, \texttt{QtUiTools} or any module whose name contains
\texttt{widgets}.

\section{Why the rule exists, twice over}

It exists for two independent reasons, and both bite.

\textbf{Testability.} There is no display on a CI runner. A Test Run has to be exercisable through
\texttt{QT_QPA_PLATFORM=offscreen}, and the suite drives real runs through the session and the engine
without ever constructing a widget. If the engine imported \texttt{QtWidgets}, every run test would
drag a GUI in, and the headless examples in \texttt{examples/harness/run_examples.py} would need a
window server.

\textbf{Thread affinity.} A Qt widget belongs to the thread that created it. The worker is a
\texttt{QThread}; the widgets live on the main thread. Code on the worker that touched a widget would
be reaching across threads into an object Qt only guarantees to touch from one thread. The engine
therefore does not know that screens exist: it emits signals, and the layout reacts.

The rule is deliberately narrow. \texttt{src/ate/report/pdf.py} imports \texttt{QApplication} inside a
function (line~350 of that file) because the PDF renderer runs a Qt WebEngine in a child process on the
main thread, and \texttt{src/ate/dialogs.py} imports \texttt{QTimer} for the same reason it needs to
arm a dialog: neither is on the worker, and neither is guarded. The guard covers the two packages that
run on the worker, and it names them in \texttt{\_GUARDED\_DIRS}.

The guard itself is tested. \texttt{test\_the\_guard\_still\_bites\_on\_a\_planted\_widget\_import}
plants each forbidden form --- including the old \texttt{PyQt6.QtWidgets} name --- and asserts the scan
reports it, and asserts that \texttt{from PySide6.QtCore import QThread, Signal} is clean. A guard
whose pattern matches nothing protects nothing, and that is the failure this test would catch.

\section{Signals are the only way across}

The engine reaches the GUI only through queued signals, and the session forwards them. The consequence
for whoever writes a new signal: the receiver runs on the thread that owns the receiver object, not on
the thread that emitted. That is the whole of Qt's cross-thread story in one sentence, and chapter
\ref{ch:engine} shows what it means for each of the engine's signals.

\cando{explain why a widget import in \texttt{src/ate/engine} is a bug rather than a style choice, and
name the two packages the guard scans.}

\chapter{The run engine: one worker, nine signals}\label{ch:engine}

\texttt{RunEngine} is a \texttt{QThread} subclass in
\texttt{src/\allowbreak{}ate/\allowbreak{}engine/\allowbreak{}run\_\allowbreak{}engine.py}. Its \texttt{run()} method is the worker: it inserts the
\texttt{test\_run} row, walks the expanded step tree, and closes the run. Everything the screen knows
about the run arrives through the nine signals the class declares.

\section{The worker's lifetime, and why it owns itself}

A \texttt{QThread} destroyed while its thread is still running makes Qt abort the process
(\texttt{QThread: Destroyed while thread is still running}: SIGABRT on Linux, 0xC0000409 on Windows,
with nothing printed). The danger is real because the session drops its engine when it is finished,
and that drop is queued while the worker is still returning from \texttt{run()}. Whichever side lets go
last then destroys a running thread.

So a started engine keeps itself alive. \texttt{RunEngine.start} adds the engine to
\texttt{\_live\_engines} and connects \texttt{finished} to \texttt{\_\allowbreak{}forget\_\allowbreak{}ended\_\allowbreak{}worker}, which
waits for the thread and then removes it from the set. \texttt{stop\_\allowbreak{}live\_\allowbreak{}engines(timeout\_\allowbreak{}ms=5000)}
aborts every worker still in that set and waits for it, and it is registered with \texttt{atexit}, so a
worker running when the interpreter exits is not destroyed running. If it returns a non-empty list, the
process is about to be killed by Qt, and the test suite turns that into an error message that explains
itself (chapter~\ref{ch:failed-run}).

Two rules follow, and both are easy to break by accident:

\begin{itemize}
  \item Do not \texttt{parent} a \texttt{RunEngine} to a widget in the hope that Qt will clean it up.
    A parented engine is still in \texttt{\_live\_engines}, and the widget can be deleted underneath a
    live worker.
  \item Do not call \texttt{engine.wait()} on the main thread without first calling
    \texttt{request\_abort()}. \texttt{stop\_and\_wait} does both: it sets the abort flag, then waits
    \texttt{timeout\_ms}. \texttt{wait()} releases the GIL while it waits, so the worker keeps working.
\end{itemize}

\section{What the worker does, in order}

\texttt{run()} is a guard. It calls \texttt{\_execute()} inside \texttt{try/except}, and an exception
that escapes is turned into \texttt{\_\allowbreak{}end\_\allowbreak{}unexpectedly(trace)}: the run is recorded as
\texttt{INTERRUPTED}/\texttt{FAIL}, the message is printed to stderr, the application log gets it
through the \texttt{on\_error} callable, every open Device is closed, and \texttt{FINISHED} is emitted
anyway. Each of those steps is attempted separately, because the exception may have come from the store
itself. Post-test Actions do not run after an unexpected end, as they would not after a crash.

This guard matters more since packet 32a: under PyQt6 an exception leaving \texttt{QThread.run} aborted
the process; under PySide6 it prints \texttt{Error calling Python override of QThread::run()} and the
thread simply ends. Without the guard, the run row would stay open and the screen would wait for a
\texttt{FINISHED} that never arrives. \texttt{tests/\allowbreak{}test\_\allowbreak{}unexpected\_\allowbreak{}run\_\allowbreak{}end.py} proves the run is
recorded \texttt{INTERRUPTED} and \texttt{FINISHED} is still announced.

\texttt{\_execute()} then does the work in this order:

\begin{enumerate}
  \item build a \texttt{LuaHost} for the schedule, which loads and validates it, expands
    Sub-schedules, and registers the Device drivers;
  \item insert the \texttt{test\_run} row and emit \texttt{run\_uuid\_reported};
  \item run every Pre-test Action, in Pre-test order, recording an \texttt{action} row for each; a
    Pre-test Action that raises ends the run \texttt{PRETEST\_FAILED} and records every step NOTRUN;
  \item walk the step list: header rows record their place in the tree and never run; disabled steps
    record NOTRUN with no Attempt; the abort flag is checked at every step boundary;
  \item for each enabled step, run its action closure once per Attempt, recording one \texttt{attempt}
    row each. A step's result is its last Attempt's; a step whose action raised records an
    \texttt{ERROR} Attempt; a step with no report records \texttt{NOREPORT};
  \item pause on a Stop-on-Fail failure until the Operator answers (section below);
  \item run the Post-test Actions, which run even after an abort (spec story 84), each ending on a second
    abort;
  \item record the Media Release chain and every file the run used, close open Devices, and write the
    run's result and End Reason.
\end{enumerate}

\section{The nine signals}

Every one of these is emitted on the worker and delivered to the main thread, where the receiver's own
event loop runs the slot. \texttt{tests/\allowbreak{}test\_\allowbreak{}signal\_\allowbreak{}threads\_\allowbreak32a.py} drives a real run with a
Stop-on-Fail pause, a POLL and a blocking input, and records the thread each slot ran on: they all ran
on the main thread. The same test connects one signal with \texttt{DirectConnection} as a control, and
sees the worker's thread --- and planting \texttt{DirectConnection} in the main test makes it fail. A
new signal has to be added to that test.

\begingroup\small
\begin{longtable}{p{3.3cm} p{3.6cm} p{5.2cm}}
\toprule
\textbf{Signal} & \textbf{Arguments} & \textbf{When it is emitted} \\
\midrule
\endhead
\texttt{run\_uuid\_reported} & the run key & as soon as the \texttt{test\_run} row exists; also emitted
  by an unexpected end, so the report can still find the row \\
\texttt{state\_changed} & \texttt{RunStatus} & \texttt{IDLE} $\to$ \texttt{RUNNING} $\to$ \texttt{FINISHED} \\
\texttt{input\_asked} & node seq, Attempt no, input index & a blocking \texttt{ATE.input(n)} has started
  waiting; announced from \texttt{\_on\_asked} while the worker sits in \texttt{get\_raw} \\
\texttt{step\_finished} & node seq & this step's Attempts are done, so its POLL inputs close \\
\texttt{failure\_prompt} & node seq, Attempt no & a Stop-on-Fail failure is waiting for a choice \\
\texttt{prompt\_answered} & node seq, Attempt no, choice & the pause is over: emitted when answered, and
  again when the worker moves on \\
\texttt{action\_started} & \texttt{"pre"}/\texttt{"post"}, completed, total & a Pre- or Post-test Action
  is about to run, so the status bar can say \texttt{Pre-test n/total} \\
\texttt{step\_started} & node seq & an Attempt is about to run, again on each Retry \\
\texttt{log\_lines\_recorded} & the run key & as soon as a batch of log rows lands in the store, so a
  layout can paint the newest lines while the worker is still running (packet 40a) \\
\bottomrule
\end{longtable}
\endgroup

\texttt{prompt\_answered} is emitted twice on purpose (packet 25b review): the screen clears its prompt
as soon as the answer is given, and a \texttt{failure\_prompt} still queued behind it must not leave an
answered prompt painted.

Two types of argument deserve a warning. \texttt{Signal(int)} in PySide6 is a C \texttt{int}:
\texttt{emit(2**31)} with a slot connected returns normally and the slot receives -2147483648, with a
\texttt{RuntimeWarning}; \texttt{Signal(str)} turns a non-string into \texttt{''}. A slot cannot see a
bad value, so the check is on the sending side, and
\texttt{tests/\allowbreak{}test\_\allowbreak{}signal\_\allowbreak{}threads\_\allowbreak32a.py} pins every value the engine sends. Everything the engine
emits is bounded by the schedule (node sequence, Attempt number, input index) or is a string that is
always a string (\texttt{"pre"}/\texttt{"post"}, \texttt{CONTINUE}/\texttt{RETRY}/\texttt{ABORT}, a
path). Anything that could exceed \texttt{2**31 - 1} must be declared \texttt{Signal(object)}.

\section{Blocking, and how a block lets go}

Three things block the worker, and each has to give way to an abort:

\begin{itemize}
  \item \textbf{A Stop-on-Fail pause.} \texttt{\_\allowbreak{}pause\_\allowbreak{}for\_\allowbreak{}operator} blocks with no timeout until the
    Operator answers. Packet 32b found the bug that this is about: a run-level abort that arrived before
    the pause was entered used to be lost, because the pause never read the flag. It now ends on the
    abort. A real answer given at the same moment is not overwritten, and \texttt{choose} returns
    \texttt{None} once the flag is set, so the order of the two decides rather than the worker's 50 ms
    poll.
  \item \textbf{\texttt{ATE.input(n)}.} A blocking input waits in \texttt{OperatorInput.get\_\allowbreak{}raw}, which
    takes the abort flag and gives up. The ask is announced with \texttt{input\_asked} first, so a layout
    can put a widget under the step and the Operator can answer while the worker sits.
  \item \textbf{\texttt{ATE.sleep(seconds)}.} The wait is sliced, so the abort lands within about 50 ms.
\end{itemize}

A busy loop is interrupted by the Lua count hook, which calls \texttt{aborted()} every 1000
instructions and raises \texttt{aborted} once the flag is set; \texttt{pcall} and \texttt{xpcall} are
wrapped to re-raise it so an abort cannot be caught. A Device call is interrupted by
\texttt{ate/devices/abort.py} (chapter~\ref{ch:devices}).

Three things cannot be interrupted at all, and the manual says so rather than promising otherwise: a
SQLite commit already in progress, one Lua C call such as a pattern match on a huge string, and
\texttt{set\_visa\_attribute} or \texttt{resource.timeout} on the worker.

\section{What may run on the worker}

Allowed, because the abort flag is checked inside it and it never touches a widget:

\begin{itemize}
  \item a Test Step or Action closure, synchronously. \texttt{run\_engine.py}'s docstring assumes
    this: the sandbox does not remove coroutines, but spec story 83 records a busy loop in one as an
    exception to the abort bound, so \texttt{ATE.report} is attributed to a step by snapshotting the
    host's report log before the closure runs and diffing it afterwards;
  \item store writes. Each is its own commit (the Engine is autocommit), which is what makes the record
    append-only and crash-safe, and also why a run with many rows costs many commits --- see
    chapter~\ref{ch:storage};
  \item Device reads, writes, opens and closes, through the registry, which hands every call to
    \texttt{run\_abortable};
  \item the drains: \texttt{\_\allowbreak{}flush\_\allowbreak{}device\_\allowbreak{}warnings}, \texttt{\_flush\_equipment} and
    \texttt{\_record\_files} run on the worker and write what the host collected.
\end{itemize}

Not allowed:

\begin{itemize}
  \item touching a widget or importing \texttt{QtWidgets};
  \item a wait that does not check the abort flag. Every loop a worker enters must have an exit the
    abort can reach;
  \item emitting from the main thread into the engine as if it were the worker. The main thread calls
    \texttt{choose()} and \texttt{request\_abort()}; it does not step the schedule.
\end{itemize}

One detail that trips people reading log stamps: \texttt{record\_log} takes \texttt{time.monotonic()}
when the line is recorded, and the drain only maps that to wall-clock. The \texttt{at} column promises
when the worker logged the line, not when the worker got there on time. Packet 31e's stamp test was
comparing it with a deadline and was measuring the wrong thing.

\cando{trace a run from \texttt{start()} to \texttt{FINISHED}, name each signal and what blocks the
worker, and say what an abort can and cannot interrupt.}

\chapter{The run session: what a layout binds to}\label{ch:session}

\texttt{HeadlessRunSession} in \texttt{src/\allowbreak{}ate/\allowbreak{}run\_\allowbreak{}session/\allowbreak{}headless.py} is the object a screen talks
to, and the object a headless run uses. It is a \texttt{QObject}: that is what lets a widget connect to
its signals. It imports no widgets, and it writes no SQL. It holds the Operator's decisions, owns one
engine at a time, and turns the worker's signals into its own.

\section{The states}

\texttt{SessionState} has four values, and the session keeps \texttt{\_state} in step with what it
emits: \texttt{IDLE} (nothing running, fields not yet entered), \texttt{SELECTING} (a schedule is
selected, the run has not started), \texttt{STARTED} (the worker is running), \texttt{FINISHED}. The
engine's \texttt{RunStatus} is mapped onto these by \texttt{\_map\_run\_status}: \texttt{RUNNING}
becomes \texttt{STARTED}.

\texttt{STARTED} is announced on the calling thread before the worker starts, so a screen can paint the
run view before the first signal arrives. The worker's queued signals then move the state.

\section{The nine signals}

\begin{itemize}
  \item \texttt{state\_changed}, \texttt{input\_asked}, \texttt{step\_finished}, \texttt{failure\_prompt},
    \texttt{prompt\_answered}, \texttt{action\_started}, \texttt{step\_started}: forwarded from the
    engine, so a layout connects to the session and never to the engine.
  \item \texttt{input\_answered} (node seq, Attempt no, index, value): the session's own, emitted when an
    answer it accepted reached the worker. Its value type is \texttt{object}, because a Step Input can be
    a number, a string or a list of labels.
  \item \texttt{report\_written} (HTML path, PDF path): fired exactly once per finished run, with an
    empty second path when there is no PDF --- because the render was skipped, or because it failed.
\end{itemize}

That last one is worth the space it takes. The report is written after \texttt{FINISHED} is announced,
and the PDF is rendered by a \texttt{QProcess} the session owns (chapter~\ref{ch:pdf}). The sandbox
question is asked \emph{after} the HTML is on disk, so a declined PDF still leaves the report and the
signal still fires once. A session that never gets a \texttt{report\_written} is a session whose teardown
was wrong.

\section{Rules the session enforces}

The session, not the screen, is the authority:

\begin{itemize}
  \item \texttt{start()} refuses to run with no schedule selected, and refuses a run with an empty
    Project (spec 08/11). A screen that disables the Start button is being helpful; the session is what
    makes it true.
  \item Before a run starts, \texttt{close\_\allowbreak{}interrupted\_\allowbreak{}runs} closes every run row with no
    \texttt{ended\_at} --- one left open by a crash or a power cut --- recording it
    \texttt{INTERRUPTED}/\texttt{FAIL}, and logs how many. That is spec rule 98, and it is why a run
    never stays open across launches.
  \item Stop on Fail is decided in one place: the Operator's explicit choice wins, otherwise the
    schedule's own default (captured when the schedule was selected, \texttt{True} when unset). Whatever
    it resolves to is what the run row records.
  \item \texttt{supply\_input} and \texttt{click\_poll} check the answer before the worker sees it: the
    step and Attempt must be the ones running, the index must be declared by that step, the input must
    actually be waiting, and the value must satisfy the declaration (packet 10a review). A bad value
    raises \texttt{ValueError} and is never recorded. \texttt{index} here is the 1-based \texttt{n} of
    \texttt{ATE.input(n)}; \texttt{OperatorInput} keys its channels from 0, and this is where the two
    meet.
  \item \texttt{pause\_pending()} is how a layout knows the run is paused. The layouts use the session,
    not their own prompt bar, because the bar of the layout that is not showing says nothing.
\end{itemize}

\section{Closing, and why the order matters}

\texttt{close()} ends the worker, writes the HTML report if it has not been written, marks the session
closed so a still-queued \texttt{FINISHED} finds nothing to write, stops the PDF helper, closes the
store, and closes the application log only if the session opened its own (a log handed in by the launch
belongs to the launch). \texttt{\_\_del\_\_} does the helper stop, because a session dropped without
\texttt{close()} leaves its \texttt{QProcess} to be destroyed by Qt mid-render, and the job's
\texttt{finished} callback then fires against a session whose Qt side is gone (packet 25b review: this
was the teardown crash).

Do not close a store while a worker is still writing to it. \texttt{close()} calls
\texttt{stop\_and\_wait()} first for exactly that reason; if the worker does not stop within 5 s it logs
an error rather than pretending it did.

\section{Driving a run with no display}

This is the shape the suite and the harness use, and it is the shape to copy for a new headless test.
The runner below is \texttt{docs/manuals/examples/developer/runner.py}, the file the manual's worked
example in chapter~\ref{ch:device-type} runs; \texttt{tests/test_manuals.py} runs it as a subprocess.

\lstinputlisting[language=ATEPython, caption={\texttt{docs/manuals/examples/developer/runner.py}}]{../examples/developer/runner.py}

There is no widget and no display, and the queued signals still deliver, because delivery rides the
event loop of the receiving thread. Without the \texttt{processEvents} pump nothing arrives: no
\texttt{state\_changed}, no step recorded, no report. A Stop-on-Fail pause blocks the worker until
somebody answers it, and in a headless run the script is the only somebody, which is why the runner
answers every pause from \texttt{--choice}. A blocking Step Input is answered the same way, from
\texttt{--input STEP NAME=VALUE}: \texttt{input\_asked} says which node is waiting, the Step Name is
looked up in the schedule's tree, and \texttt{supply\_input} hands the value over. The script exits
0 only when the session reached \texttt{FINISHED}.

\cando{start a Test Run from a script, answer its Step Inputs, abort it, and read the record back from
the store.}

\chapter{The screens: layouts, shared widgets, and dialogs that answer themselves}\label{ch:ui}

GUI code lives only in \texttt{src/ate/ui/} and \texttt{src/ate/app.py}. A layout contains no test
logic: it connects to the session, paints, and sends the Operator's clicks back. That separation is what
makes a Test Run testable with no display, and it is what the dialog rules below protect.

\section{Two layouts, one set of widgets}

\texttt{ui.layout} selects \texttt{staged} (the default, \texttt{StagedRunScreen} in
\texttt{src/\allowbreak{}ate/\allowbreak{}ui/\allowbreak{}staged\_\allowbreak{}run\_\allowbreak{}screen.py}) or \texttt{threePane} (\texttt{ThreePaneRunScreen} in
\texttt{src/\allowbreak{}ate/\allowbreak{}ui/\allowbreak{}three\_\allowbreak{}pane\_\allowbreak{}run\_\allowbreak{}screen.py}). \texttt{App.set\_layout} swaps the two in a
\texttt{QStackedWidget} --- only if \texttt{ui.allowLayoutSwitch} is on, and for that session only, since
the config is never written by a View switch. \texttt{\_carry\_over} moves the Operator's state --- the
four setup fields, the filter, the On-failure pair, the picked schedule --- from the layout being left to
the one being shown, so switching mid-session does not lose a run setup. The layout that is no longer
showing is made inert with \texttt{set\_active(False)}, and the one coming to the front takes the keyboard
through \texttt{focus\_waiting()}. \texttt{\_mark\_layout} stamps the window so a test can tell which one
is current.

The pieces they share are in \texttt{src/\allowbreak{}ate/\allowbreak{}ui/\allowbreak{}run\_\allowbreak{}widgets.py}:

\begin{itemize}
  \item \texttt{LogPane} --- the Run log, with the level/source/owner formatting both layouts use;
  \item \texttt{StatusBar} --- the run's phase and the current-item timer (packet 14d);
  \item \texttt{StepEnabler} --- the enable tree: a header's check state is derived from its subtree
    (packet 15d), and \texttt{toggle} returns the rows it changed;
  \item \texttt{ResultCounts}, \texttt{RunFieldsForm}, \texttt{SchedulePicker} (which emits
    \texttt{selection\_changed} with the exact schedule name, or \texttt{""} when nothing runnable is
    selected), \texttt{FailurePromptBar}, \texttt{StepInputWidget}, \texttt{RunScreenMixin} and
    \texttt{MismatchConfirmation}.
  \item \texttt{load\_\allowbreak{}schedule\_\allowbreak{}steps(session, name)} is how a layout gets the expanded step list before
    a run, and \texttt{\_\allowbreak{}broken\_\allowbreak{}row\_\allowbreak{}widget} is how a schedule that fails to load is shown with its error.
\end{itemize}

The other screens: \texttt{HistoryScreen} (\texttt{src/\allowbreak{}ate/\allowbreak{}ui/\allowbreak{}history\_\allowbreak{}screen.py}, ticket 23b) and
\texttt{ExportLogDialog} (\texttt{src/\allowbreak{}ate/\allowbreak{}ui/\allowbreak{}export\_\allowbreak{}log\_\allowbreak{}dialog.py}, packet 24b). History queries go
through \texttt{ate/\allowbreak{}history/\allowbreak{}session.py}, which has no widgets and is testable on its own.

New widgets go in \texttt{run\_widgets.py} if both layouts need them. A widget in one layout that the
other layout also ends up needing is how the two drift apart, and the layout-switching tests
(\texttt{tests/\allowbreak{}test\_\allowbreak{}layout\_\allowbreak{}switching\_\allowbreak25b.py}) are what catch the state loss.

\section{Every modal must be able to answer itself}

\texttt{src/ate/dialogs.py} is the only place that decides when a dialog has nobody to answer it. A modal
\texttt{exec()} blocks until someone clicks. On a CI runner, in a container, or in any process started
with \texttt{QT\_\allowbreak{}QPA\_\allowbreak{}PLATFORM=offscreen} there is no window server, so no one ever clicks: the second real
CI run hung for 30 minutes on the run-start mismatch confirmation, and the first hung on the
unsandboxed-PDF question.

\texttt{dismiss\_\allowbreak{}if\_\allowbreak{}headless(dialog, answer)} arms a dialog to close itself as soon as its own event
loop runs, using \texttt{QTimer.singleShot(0, ...)}. It fires when either of these is true:

\begin{itemize}
  \item \texttt{QT\_QPA\_PLATFORM} is \texttt{offscreen} or \texttt{minimal} --- no display at all;
  \item \texttt{ATE\_\allowbreak{}AUTO\_\allowbreak{}DISMISS\_\allowbreak{}DIALOG} is set to something that means yes (\texttt{""}, \texttt{"0"},
    \texttt{"false"}, \texttt{"no"} and \texttt{"off"} all mean unset). That variable is the explicit
    seam for subprocess tests (packet 30c); it does not make a real operator's dialog auto-answer on a
    desktop.
\end{itemize}

With no \texttt{answer}, the box is \texttt{accept()}ed, which clicks no button:
\texttt{clickedButton()} is \texttt{None}, and every prompt therefore takes its \emph{negative} branch ---
the missing-config prompt aborts, the close confirmation keeps the run going, the PDF question is answered
``Skip PDF''. A dialog with a decision of its own is handed an explicit function:
\texttt{MismatchConfirmation.answer\_\allowbreak{}from\_\allowbreak{}setting} starts the run and stores the mismatched bytes only
when \texttt{database.storeMismatchedFiles} says \texttt{always}, because an unattended run makes no
choice and a ticked-by-default box is not one (packet 31d).

The rule for a new modal is short: if you add a dialog, call \texttt{dismiss\_\allowbreak{}if\_\allowbreak{}headless} on it, and
decide what the negative branch means. \texttt{tests/\allowbreak{}test\_\allowbreak{}modals\_\allowbreak{}never\_\allowbreak{}block\_\allowbreak{}headless.py} is where
that is pinned; a dialog that is not armed hangs the suite, and the hang looks like a stuck run rather
than a missing call.

\section{What the shell does}

\texttt{app.py} reads the config once at launch (spec 17), refuses to open if the config or the database
is wrong, builds the layout named by \texttt{ui.layout}, installs the exception hook
(chapter~\ref{ch:config}), and shows the History screen when \texttt{ui.allowHistoryDuringRun} allows
it. \texttt{launch()} is the entry point, and \texttt{tests/\allowbreak{}test\_\allowbreak{}gui\_\allowbreak{}launch\_\allowbreak30a.py} drives it with the
dialogs injected, so no window is ever shown in the suite.

\cando{add a modal without hanging CI, and say which layout a widget belongs in.}

\chapter{Storage: the schema, batching, and the log counter}\label{ch:storage}

\texttt{src/ate/storage/} is the only place ATE writes SQL. Three files: \texttt{schema.py} declares
every table as a SQLAlchemy Core table, \texttt{store.py} implements \texttt{RunStore}, the only
interface the rest of the application uses, and \texttt{backup.py} holds the one routine that backs a
database up and then migrates it. Nothing outside these files contains a table name or a query.

\section{The record, and the version}

The Test Run record is ten tables. \texttt{tests/\allowbreak{}test\_\allowbreak{}manuals.py} reads them from
\texttt{schema.metadata.tables}, so this list is the code's list:

\begingroup\small
\begin{longtable}{p{3.2cm} p{8.4cm}}
\toprule
\textbf{Table} & \textbf{One row is} \\
\midrule
\endhead
\texttt{test\_run} & one Test Run: Operator, Project, UUT, schedule, Media Release, Stop on Fail, start
  and end times, UTC offset, Result, End Reason \\
\texttt{node} & one step or header in the expanded tree, with its number, enable state, Result and
  Attempt count \\
\texttt{attempt} & one execution of one step, with its times, Result, message, actual, expected, and the
  Stop-on-Fail choice made at it \\
\texttt{action} & one Pre-test or Post-test Action call, with phase, order, times and outcome
  (\texttt{ok}, \texttt{error}, \texttt{aborted}, \texttt{skipped}) \\
\texttt{log\_line} & one line of the Run log, owned by an Action or an Attempt, or run-level \\
\texttt{input\_value} & one Step Input answer \\
\texttt{equipment} & one \texttt{ATE.addEquipment} declaration, in declaration order \\
\texttt{media\_release} & one Media Release in the resolved chain; the release owning the run is seq 0 \\
\texttt{run\_file} & one file the run used, with expected and computed checksums and whether they differed \\
\texttt{schedule\_instance} & the top-level schedule and every Sub-schedule include, in tree order \\
\bottomrule
\end{longtable}
\endgroup

Types are deliberately portable: UUIDs are 36-character text, times are ISO-8601 UTC text with
milliseconds (\texttt{\_now\_iso}), booleans are 0/1 (\texttt{\_as\_int}), structured values are JSON
text. There is no SQLite-only column type anywhere; a Postgres backend is meant to be a connection
change. The version lives in SQLite's \texttt{PRAGMA user\_version}, and \texttt{SCHEMA\_VERSION} is
currently 5.

At launch, \texttt{check\_\allowbreak{}database\_\allowbreak{}at\_\allowbreak{}launch} inspects the file read-only --- it never creates it ---
and compares \texttt{user\_version} with \texttt{SCHEMA\_VERSION}. A missing, empty or older database
refuses with \texttt{run `ate db migrate` to create or upgrade the database at \dots}; a newer one refuses
with \texttt{this ATE is too old for the database at \dots; update ATE}. Only an exact match allows
launch. The application never migrates; \texttt{ate db migrate} does, and the hint carries
\texttt{--config <path>} when the Station is not reading the platform default (packet 30c).

\texttt{RunStore.\_\allowbreak{}ensure\_\allowbreak{}schema} is idempotent: \texttt{metadata.create\_all} reflects each table through
the backend's own catalog, so reopening a populated database never fails. The version stamp is what
packet 35 changed: \texttt{PRAGMA user\_version} is raised only when the build's
\texttt{missing\_tables(path)} was non-empty --- when a table was really created. A database that stores
an older version but already holds every table this ATE declares has had no migration to earn the stamp,
so opening it raises \texttt{SchemaVersionError} naming that version instead of quietly declaring it
current. WAL and the version are the two pragmas set on the store's own stdlib connection: the only raw
SQL in ATE, and it is confined to pragmas.

\texttt{src/ate/storage/backup.py} is the one routine \texttt{ate db migrate} and \texttt{ate setup}
both call, so an upgrade through an installer and an upgrade by hand take the same copy.
\texttt{backup\_for\_migration} decides honestly whether a copy is needed: none for a database that does
not exist, none for one with no tables, none for one that needs no migration; otherwise SQLite's
\texttt{Connection.backup} writes \texttt{<db>.backup-<YYYYMMDD-HHMMSS>} beside the file, never over an
existing name. The backup API is used rather than \texttt{cp} because it folds \texttt{-wal} into the
copy: a database with a Test Run in flight is copied whole. \texttt{migrate\_database} then refuses if
the copy cannot be written, and reports the copy's path when a migration fails after it.

\section{One commit per row, and why that is a design decision}

The Core Engine is autocommit: each write commits immediately. That is what makes the record
append-only and crash-safe --- an open run simply has no \texttt{ended\_at} yet --- and it is also the
reason an abort over 400 unrun steps used to make 404 commits and take longer than spec story 83's 10 s with
no stall at all (packet 32b). The answer was batching, not a different transaction model.

SQLite limits the number of variables one statement may carry; the lowest default is 999
(\texttt{\_\allowbreak{}SQL\_\allowbreak{}MAX\_\allowbreak{}VARIABLES}, the pre-3.32 default, so the code stays safe on old builds too).
\texttt{\_\allowbreak{}rows\_\allowbreak{}per\_\allowbreak{}statement(table)} computes the batch size from the table's column count ---
\texttt{999 /\allowbreak/\allowbreak len(table.columns)} --- rather than a hard-coded 80 chosen for \texttt{node}'s twelve
columns. A batch sized for one table breaks the day another table gains a column, which
\texttt{test\_\allowbreak{}a\_\allowbreak{}row\_\allowbreak{}batch\_\allowbreak{}is\_\allowbreak{}sized\_\allowbreak{}to\_\allowbreak{}the\_\allowbreak{}table\_\allowbreak{}it\_\allowbreak{}writes} is there to prevent.
\texttt{\_row\_chunks} then yields the rows in pieces that fit.

There is a trap here that a new writer will hit. Each chunk is its own commit, so the run's counters must
advance \emph{inside} the loop, only for chunks that landed. \texttt{insert\_log\_lines} does exactly
that: before packet 32b a half-written batch left \texttt{\_log\_seq} behind the rows already in the
table, and the next write failed with \texttt{UNIQUE constraint failed: log\_line.run\_uuid,
log\_line.seq}. The test is
\texttt{test\_\allowbreak{}a\_\allowbreak{}log\_\allowbreak{}batch\_\allowbreak{}failing\_\allowbreak{}part\_\allowbreak{}way\_\allowbreak{}through\_\allowbreak{}keeps\_\allowbreak{}the\_\allowbreak{}log\_\allowbreak{}counter\_\allowbreak{}true}.

\section{\texttt{\_log\_seq}: one counter per run, four tables}

\texttt{log\_line}, \texttt{action}, \texttt{input\_value} and \texttt{equipment} all have
\texttt{(run\_uuid, seq)} as their primary key, and they share one per-run counter held in
\texttt{RunStore.\_log\_seq}. That is what makes a log line, an Action and an Equipment declaration
uniquely keyed without a global sequence, and it is why the ordering of the drain calls matters: the
engine diffs the host's collected lines around each step, and a batch can hold several lines from one
Attempt.

Rules for a new row type that shares this key:

\begin{itemize}
  \item take \texttt{next\_\allowbreak{}seq = self.\_\allowbreak{}log\_\allowbreak{}seq.get(run\_\allowbreak{}uuid, 0)} and advance it by the number of rows
    that actually landed;
  \item never assume a batch is one commit;
  \item stamp ownership (\texttt{action\_seq}, or \texttt{node\_seq} plus \texttt{attempt\_no}) the way the
    log lines do, so the report can group a row under the step or Action that produced it.
\end{itemize}

\section{Reading, and the \texttt{\_run} subtlety}

Every read goes through \texttt{RunStore.\_run}, which executes a Core statement and
\texttt{fetchall()}s the rows \emph{inside} the connection context. Packet 25a review found why:
returning the open \texttt{Result} left a read statement alive after the context closed; the GUI thread
fetched it later, the connection pool handed the same sqlite3 connection --- cursor still open --- to the
run worker, and a worker INSERT colliding with the unfinalised read aborted a long run with ``database is
locked''. Do not return a result object from \texttt{\_run}. Materialise it.

\texttt{query\_runs} is the History screen's query: filters on the eight \texttt{HISTORY\_COLUMNS},
sorting on exactly those keys, and \texttt{run\_uuid} ordered last so paging never repeats or skips a row
when several runs share a sort value. Dates are normalised to ISO-8601 UTC by \texttt{\_norm\_date} so
lexical comparison is correct.

\cando{add a table to the record the way the code does it, and explain why a batch is not one commit.}

\chapter{The Lua host and the sandbox}\label{ch:lua}

\texttt{src/ate/engine/lua\_host.py} is 1798 lines, and most of them are about one thing: a Test
Schedule is untrusted input, and the host has to run it without letting it reach the machine. Read it as
three parts: the runtime and its lockdown, the \texttt{ATE} table, and the schedule builders that turn a
Lua file into the tree the engine walks.

\section{One runtime per schedule load}

Each load makes a fresh \texttt{lupa} runtime. Nothing is shared between schedules in a run, so a
schedule cannot read another schedule's globals, and a Test Run's state does not leak from one step to
another through the interpreter. The runtime is created with \texttt{register\_builtins=False}, so
Python's builtins are not present at all, and then \texttt{\_lockdown} shadows the file-loading entry
points:

\begin{itemize}
  \item \texttt{dofile}, \texttt{loadfile} and \texttt{load} are set to \texttt{nil}. The Lua standard
    library still provides them, so they must be shadowed explicitly; calling a nil value raises
    ``attempt to call a nil value'', which the engine turns into an ERROR Attempt.
  \item \texttt{python} is set to nil, so the host's Python is not reachable from a schedule.
  \item \texttt{require} is replaced by ATE's own, and \texttt{package.path} and \texttt{package.cpath}
    are emptied (packet 06a), so only the whitelist from the Test Definition --- and the modules bundled
    in \texttt{src/ate/lua/}, reachable as \texttt{require("ate.devices")} --- can be loaded.
  \item \texttt{debug} is reduced to \texttt{debug.traceback}. With the full library a schedule could
    reach the instance loader through \texttt{debug.getregistry}, strip the read-only locks with
    \texttt{debug.setmetatable}, read shared globals past \texttt{\_\_metatable\_\_} or remove the abort
    hook. ATE keeps what it needs on the Python side; only \texttt{debug.traceback} stays, because error
    tracebacks are built through it.
  \item Each schedule instance --- the top level and every include --- gets its own global table, built by
    a loader that captured the stock \texttt{load} before the lockdown removed it.
\end{itemize}

The lockdown runs \emph{before} anything else is registered, so a schedule cannot escape at load time
either. \texttt{tests/test\_lua\_isolation.py} and \texttt{test\_lua\_lockdown\_calls.py} are the tests
that prove this, and the example \texttt{forbidden.tslua} in
\texttt{docs/manuals/examples/test-writer/} is a schedule that opens a Device while the file is loading,
which the host refuses.

The abort hook is installed before the schedule runs: \texttt{debug.sethook(check\_abort, "", 1000)}, and
\texttt{pcall}/\texttt{xpcall} are wrapped to re-raise \texttt{aborted} so a schedule's own error handling
cannot swallow an abort (packet 13c). \texttt{print} is redirected to \texttt{ATE.log(INFO, \dots)}, so
schedule output is always logged.

\section{The \texttt{ATE} table}

\texttt{\_AteTable} is what a schedule sees. Two things about it are not obvious:

\begin{itemize}
  \item every function is wrapped by \texttt{\_isolate}, which turns a raised ATE error into the
    ``(false, message)`` Device failure shape or an ERROR Attempt, and names the schedule file and line
    in the message;
  \item \texttt{ATE.run} is a read-only Lua table built from \texttt{\_RUN\_FIELDS}:
    \texttt{UUT\_serial}, \texttt{Operator}, \texttt{Station}, \texttt{Project},
    \texttt{Test\_Schedule\_Name}, \texttt{Media\_Release\_Media\_Number}, \texttt{name}, \texttt{version}.
    Writing to it raises \texttt{ATE.run: read-only}. It is rebuilt by \texttt{publish\_run} through
    \texttt{rawset}, which is how the engine fills it without tripping the guard.
\end{itemize}

The enums a schedule can use are \texttt{Result}, \texttt{Log}, \texttt{DeviceType}, \texttt{Parity} and
\texttt{FlowControl}, declared by \texttt{\_declare\_enums} as read-only tables, and
\texttt{tests/test_manuals.py} derives them from \texttt{\_ENUMS} so a manual cannot list fewer than the
code has.

Context matters: \texttt{\_context} is one of \texttt{"load"}, \texttt{"pretest"}, \texttt{"posttest"} or
\texttt{"step"}, and the functions that are only legal in one of them raise when called in another. A
Device opened at load time is a load error, not a runtime error.

\section{Adding an \texttt{ATE} function}

A new function has to do four things or it is not really available:

\begin{enumerate}
  \item define it on \texttt{\_AteTable} and \texttt{register} it, so the table exposes it;
  \item make it record what it did through the host (the \texttt{logs}, \texttt{reports} or
    \texttt{equipment} lists) rather than writing to the store directly --- the engine owns the rows and
    attributes them to the step or Action running;
  \item check the context if it is only legal in one phase;
  \item add it to both manuals. \texttt{ate\_table\_functions()} in
    \texttt{tests/test\_manuals.py} reads the table off the object, so a function the code has and the
    manual does not name fails the suite.
\end{enumerate}

Seams excluded from that list are \texttt{register} and \texttt{publish\_run}, which build the table
rather than being calls a schedule makes.

\cando{say why a schedule cannot open a file, and add a function to the sandbox without breaking the
coverage test.}

\chapter{Device drivers, the registry, and how an abort reaches a call}\label{ch:devices}

A Device is anything a Test Schedule opens: a TCP socket, a UDP peer, a serial port, a VISA resource.
Each Device Type has one driver, and the driver contract is small on purpose --- a new Device Type is a
file and one line, and chapter~\ref{ch:device-type} walks through it with a driver that the test suite
actually runs.

\section{The contract}

\texttt{ate/\allowbreak{}devices/\allowbreak{}registry.py} defines \texttt{Driver}, an abstract base class with three class
attributes and two methods. Only \texttt{open} is abstract. Teardown is not the driver's job: the handle
\texttt{open} returns has its own \texttt{close}, and that is what the registry calls. \texttt{Driver.close}
is declared with a no-op default, but nothing reaches it --- none of the four shipped drivers is
teared down through the driver, and \texttt{VisaDriver.close} is the only override, dead code kept so the
contract reads as written.

\begin{lstlisting}[language=ATEPython]
class Driver(ABC):
    name: ClassVar[str] = "driver"                       # the Device Type name, e.g. "TCP"
    accepts_receive_buffer_bytes: ClassVar[bool] = False # does open() take the buffer cap?
    fields: ClassVar[dict[str, tuple[tuple[type, ...], bool]]] = {}  # field -> (types, required)

    @abstractmethod
    def open(self, params: dict[str, Any]) -> Any: ...

    def close(self, device: Any) -> None:
        """Release a device. The registry closes the handle, never through this method."""
\end{lstlisting}

\texttt{validate\_params} checks every declared field generically: an unknown field, a missing required
one, or a value of the wrong type raises a \texttt{ValueError} naming the field. Booleans are explicitly
not numbers --- \texttt{isinstance(True, int)} is true in Python, so a driver that accepts
\texttt{(int, float)} must reject a boolean, and \texttt{\_matches} does that.

The registry never knows how a driver connects. It only knows that \texttt{open} produced a handle, and
it tracks that handle until \texttt{close\_device} is called on it, so a Device left open at the end of a
run can be closed and warned about (packet 16a). \texttt{close\_device} is the whole of the teardown
contract: it calls the handle's own \texttt{close}, through \texttt{run\_abortable} when the close is one
the abort is meant to end, and \texttt{ATE.close} and the run's end-of-run cleanup both go through it.
\texttt{\_StreamDevice.close} (\texttt{stream.py}) stops the reader thread, closes the transport and joins
the reader; \texttt{\_VisaDevice.close} (\texttt{visa.py}) closes the resource. Both are safe to call
twice, so a Device closed by the schedule and then closed again by the run's cleanup is not an error.

The four drivers that ship, and the Device Types they serve, are \texttt{TcpDriver} (TCP),
\texttt{SerialDriver} (SERIAL), \texttt{UdpDriver} (UDP) and \texttt{VisaDriver} (VISA).
\texttt{tests/\allowbreak{}test\_\allowbreak{}manuals.py} derives that list by importing every module in \texttt{ate.devices} and
collecting the classes that subclass \texttt{Driver}, so a driver added to the code and not named in the
manual fails the suite.

TCP, SERIAL and UDP share \texttt{devices/stream.py}: a background reader thread, a receive buffer
capped by \texttt{devices.receiveBufferBytes} (default 16777216 bytes), an overflow that drops the oldest
bytes and reports \texttt{overflow = true} with one WARNING on the next read
(\texttt{receive buffer overflow - dropped 4096 oldest byte(s)}), and the failure and warning logic.
UDP is that same machinery with one change: \texttt{\_UdpDevice} subclasses \texttt{\_StreamDevice} and
replaces \texttt{\_reader\_loop} (\texttt{udp.py}, "The shared drain loop, except that an empty datagram
is data"), because the base loop skips a falsy chunk and a zero-length datagram is a real message. VISA
is not on the stream at all: \texttt{\_VisaDevice} in \texttt{visa.py} is its own class, reads on demand
and has no reader thread, so \texttt{accepts\_\allowbreak{}receive\_\allowbreak{}buffer\_\allowbreak{}bytes} is \texttt{False} for it --- there
is no buffer to cap.

\section{Tracing}

When \texttt{deviceTrace.enabled} is on, the registry wraps every handle it opens in
\texttt{TracedDevice} and logs each open, read, write and close as one \texttt{DEBUG} line with source
\texttt{device}, through the host's log sink --- so a trace line is owned by the step or Action that made
the call, exactly like a schedule log line. Payloads are escaped and truncated at
\texttt{deviceTrace.maxBytes}, default \texttt{DEFAULT\_\allowbreak{}TRACE\_\allowbreak{}MAX\_\allowbreak{}BYTES = 256}. A call the abort left
behind gets its own line, \texttt{open <device> aborted (ms)} or \texttt{close <device> aborted (ms)}
(packet 32b review round: before the fix the trace list was empty).

\section{Abort, and what an abort cannot stop}

Story 83 names two \emph{allowed} exceptions to the ten-second bound, and the suite measures both
rather than trying to fix them: a command a step runs with \texttt{os.execute}, which the abort
cannot end, and a busy loop inside a coroutine, where the count hook never ticks
(\texttt{tests/test\_\allowbreak{}manuals.py}: \texttt{test\_\allowbreak{}a\_\allowbreak{}command\_\allowbreak{}in\_\allowbreak{}a\_\allowbreak{}step\_\allowbreak{}holds\_\allowbreak{}the\_\allowbreak{}abort\_\allowbreak{}until\_\allowbreak{}it\_\allowbreak{}finishes},
\texttt{test\_\allowbreak{}a\_\allowbreak{}coroutine\_\allowbreak{}loop\_\allowbreak{}runs\_\allowbreak{}to\_\allowbreak{}completion\_\allowbreak{}after\_\allowbreak{}an\_\allowbreak{}abort}).
Neither is forbidden in a schedule. \texttt{os.exit} is not one of them: it ends the process with the
run row left open, and the next launch closes it as \texttt{INTERRUPTED}/\texttt{FAIL}.

\texttt{ate/devices/abort.py} is packet 32b's answer to ``the abort does not stop the worker in 10 s''.
There are two kinds of Device wait:

\begin{itemize}
  \item \textbf{A wait ATE owns} --- a buffer read, a \texttt{select}, a pyvisa-py SOCKET read or write.
    It runs in slices of \texttt{POLL\_SECONDS} ($0.05$ s) and checks the flag between them, so
    it ends in about 50 ms. \texttt{wait\_slice} caps the next slice by the call's own remaining timeout,
    and \texttt{check\_abort} raises \texttt{DeviceAborted} which the ATE wrappers turn into Lua's
    \texttt{aborted} error.
  \item \textbf{A wait only the OS or a C library can end} --- a connect, a DNS lookup, \texttt{viRead} on
    a real VISA stack, a pyserial write, a close on a stuck tty. These run on their own daemon thread
    through \texttt{run\_abortable}. The worker watches the flag and raises, leaving the call behind; it
    finishes when its own timeout says so, and its result is discarded rather than recorded.
\end{itemize}

Two details that were found by test and are easy to break again. \texttt{\_Call.\_run} catches
\texttt{BaseException}, not \texttt{Exception}, because a helper thread killed by \texttt{KeyboardInterrupt}
or \texttt{SystemExit} would otherwise leave no outcome and the worker would wait forever
(\texttt{test\_\allowbreak{}a\_\allowbreak{}helper\_\allowbreak{}thread\_\allowbreak{}dying\_\allowbreak{}on\_\allowbreak{}a\_\allowbreak{}base\_\allowbreak{}exception\_\allowbreak{}reaches\_\allowbreak{}the\_\allowbreak{}worker}). And a helper thread
left behind is joined in slices, so an aborted run's leftover closes are bounded:
\texttt{\_\allowbreak{}ABORT\_\allowbreak{}CLOSE\_\allowbreak{}SECONDS = 3.0} in \texttt{run\_engine.py} is the total for all closes at the end of
an aborted run, each with a WARNING naming the Device.

The consequence for a Station, stated in the admin manual as well: an aborted run can leave a Device
close running for seconds, a Post-test Action on the same Device waits for it, and a new run opening the
same instrument meanwhile may find it busy. That is not a bug to fix; it is the boundary of what an abort
can do.

\section{Which threads a Device runs on}

The background readers are \texttt{threading.Thread} objects, not \texttt{QThread}s. That matters twice:
their exceptions go to \texttt{threading.excepthook}, not \texttt{sys.excepthook}, so the
exception-log hook does not see them --- they catch their own --- and the architecture guard does not
scan them because they are not under \texttt{src/ate/engine} or \texttt{src/ate/run\_session}.

\cando{write a driver that the registry will accept, and explain what happens to an open Device when a run
is aborted.}

\chapter{Reports, and the PDF helper that has to be sandboxed}\label{ch:pdf}

Two files produce a Test Report. \texttt{src/\allowbreak{}ate/\allowbreak{}report/\allowbreak{}html.py} renders the self-contained HTML from
the recorded run; \texttt{src/\allowbreak{}ate/\allowbreak{}report/\allowbreak{}pdf.py} turns that HTML into PDF in a separate process, and
\texttt{src/\allowbreak{}ate/\allowbreak{}report/\allowbreak{}pdf\_\allowbreak{}job.py} is the \texttt{QProcess} plumbing that runs it. A Test Report is not
optional (spec 21a), so a failure here must never crash a run.

\section{The HTML}

\texttt{write\_\allowbreak{}report(store, run\_\allowbreak{}uuid, config)} renders through the Jinja template
\texttt{src/\allowbreak{}ate/\allowbreak{}report/\allowbreak{}templates/\allowbreak{}report.html.j2} and writes under \texttt{reports.directory} as

\begin{lstlisting}[language=ATEShell]
<UUT Type>/<UUT Type>_<serial>_<local start yyyymmdd-HHMMSS>_<PASS|FAIL>.html
\end{lstlisting}

Windows-invalid characters are replaced, and an existing file is never overwritten: a free
\texttt{-2}, \texttt{-3} suffix is used instead. The name is only free when neither the HTML nor its PDF
exists, because the PDF takes the same stem. \texttt{generated\_at} freezes the footer stamp for a
byte-deterministic write, which is what makes the report tests able to compare output.

\texttt{src/\allowbreak{}ate/\allowbreak{}report/\allowbreak{}log\_\allowbreak{}export.py} is the Run log export (packet 24a): CSV, JSON Lines or text,
filtered by level, source and owner --- the same filters \texttt{ExportLogDialog} shows.

\section{Why the PDF is a separate process}

Rendering needs Qt WebEngine. Running it in the application's process would put Chromium's renderer in
the same process as the Test Run, and the render is slow: a run's \texttt{FINISHED} must not wait for it.
So the helper is \texttt{python -m ate \_\allowbreak{}pdf <html> <pdf> --paper A4|Letter [--unsandboxed]} started by
\texttt{start\_pdf\_job}, with \texttt{PDF\_\allowbreak{}TIMEOUT\_\allowbreak{}MS = 60000}. The helper is parented to the session
that started it, and \texttt{stop\_pdf\_jobs()} must run before the session dies: a \texttt{QProcess}
destroyed mid-render by Qt makes the job's \texttt{finished} callback fire against a session whose Qt
side is gone, and that aborted the test run at teardown (packet 25b review). It is why
\texttt{HeadlessRunSession.close} and \texttt{\_\_del\_\_} both call it, and why
\texttt{conftest.py} has an autouse fixture that calls it after every test.

\section{The sandbox decision, asked once}

Chromium's sandbox cannot run as root or as a Windows administrator, and turning it off silently is a
security decision nobody should make by accident. \texttt{pdf\_\allowbreak{}sandbox\_\allowbreak{}choice()} is built once at launch
and handed to the session and to History, so neither of them reads the environment or guesses at
privilege. The decision, per PDF:

\begin{itemize}
  \item not privileged --- \texttt{PDF\_SANDBOXED}, nothing is asked. \texttt{--allow-unsandboxed-pdf}
    has no effect, and the log says so;
  \item privileged and \texttt{--allow-unsandboxed-pdf} --- \texttt{PDF\_UNSANDBOXED}, no dialog, logged
    as a WARNING;
  \item privileged and there is a dialog to ask --- the Operator's answer. No means \texttt{PDF\_SKIPPED},
    and the HTML Test Report is still written;
  \item privileged with no way to ask (headless, \texttt{ate \_selftest --run}) --- \texttt{PDF\_SKIPPED}.
\end{itemize}

Two environment variables can switch WebEngine's sandbox off by inheritance,
\texttt{SANDBOX\_\allowbreak{}OFF\_\allowbreak{}VARIABLES} --- \texttt{QTWEBENGINE\_\allowbreak{}DISABLE\_\allowbreak{}SANDBOX} and
\texttt{QTWEBENGINE\_\allowbreak{}CHROMIUM\_\allowbreak{}FLAGS} with \texttt{--no-sandbox} --- and \texttt{\_\allowbreak{}strip\_\allowbreak{}inherited\_\allowbreak{}sandbox\_\allowbreak{}off}
removes them from the child's environment so a Station cannot get an unsandboxed render without choosing
it.

The test seam is \texttt{ATE\_ELEVATED} (\texttt{ELEVATION\_ENV}): \texttt{conftest.py} pins it to 0 for
every test, because a GitHub Windows runner is an administrator and a container runs as root, and a
privileged session asks a question that \texttt{exec()} would block forever on an offscreen platform. A
test that wants to be privileged sets \texttt{ATE\_ELEVATED=1}, and it reaches a spawned \texttt{ate \_pdf}
child through the inherited environment.

\section{Proving a report rendered}

A 13-byte stub PDF is enough for signal plumbing and is what \texttt{conftest.py} installs for every
run that is not specifically about the PDF. Four tests use the real helper, because ``the PDF has text in
it'' can only be proven against a PDF that was printed: \texttt{test\_\allowbreak{}run\_\allowbreak{}end\_\allowbreak{}writes\_\allowbreak{}real\_\allowbreak{}pdf\_\allowbreak{}report},
the two \texttt{\_selftest --run} cases and the ordinary-user-with-the-flag case.
\texttt{pdf\_has\_text(pdf)} is the check. If you add a report feature, decide which side of that line
your test is on.

\cando{add something to the Test Report and know whether your test needs a real render or a stub.}

\chapter{Config, and the exception nothing caught}\label{ch:config}

\section{How the config is read}

\texttt{src/ate/config.py} is the whole of ATE's configuration. Four things a developer has to know
before touching it:

\begin{itemize}
  \item \texttt{\_SCHEMA} is the list of accepted keys, grouped by section, each with a type;
    \texttt{\_DEFAULTS} is the default for every one of them. A key the file does not have takes its
    default; a key that is not in the schema is refused with \texttt{unknown config key '<key>' in
    <file>}. \texttt{format} is required and must be 1.
  \item \texttt{<data>} is the only placeholder. It expands to the platform data directory
    (\texttt{/var/lib/ate} on Linux, \texttt{\%PROGRAMDATA\%\textbackslash ATE} on Windows). Any other
    \texttt{<word>} is a config error naming the file and the key, so a typo never becomes a directory
    literally called \texttt{<whatever>}.
  \item Config values get no environment expansion. The platform default path is decided by
    \texttt{\_platform\_default()}, and \texttt{--config <path>} beats it. Note the trap the file records:
    \texttt{sys.platform} is \texttt{win32}, never \texttt{windows}, and \texttt{\_is\_windows} is the only
    platform test used anywhere.
  \item \texttt{ATE\_DEFAULT\_CONFIG} and \texttt{ATE\_DATA\_DIR} are test seams. They are read nowhere
    else in ATE: they exist so a test can exercise ``there is no config at the default location'' in a
    real subprocess without writing to \texttt{/etc/ate} (packet 30c). Do not use them in production code.
\end{itemize}

Adding a key means adding it to \texttt{\_SCHEMA}, to \texttt{\_DEFAULTS}, to the accessor on
\texttt{Config}, and to both manuals --- \texttt{config\_keys()} in \texttt{tests/test\_manuals.py}
derives the set from \texttt{\_SCHEMA} and fails a manual that does not name it.

\section{The application log}

\texttt{src/ate/log.py} opens one \texttt{AppLog} per launch, at
\texttt{<data>/logs/ate-app-<launch time>.log}, honouring \texttt{log.level}. A session given a config
opens its own if the launch did not hand one in; a log handed in belongs to the launch and the session
does not close it. A log directory that cannot be opened must not stop the application: the failure leaves
\texttt{on\_log} as \texttt{None}, and the session falls back to printing.

\section{\texttt{UnhandledExceptionLog}, and what it does not see}\label{sec:exception}

PySide6 does not end the process when a slot, or a \texttt{QThread.run} override, raises. It hands the
exception to \texttt{sys.excepthook} and carries on --- so a worker keeps recording while a screen quietly
stops updating. \texttt{launch()} installs \texttt{UnhandledExceptionLog} on \texttt{sys.excepthook}: it
writes the traceback to the per-launch application log at ERROR, passes the exception to the previous hook
(so stderr and pytest-qt's capture still see it), and emits \texttt{reported} with the log path so the
window can show one line in its status bar:

\begin{lstlisting}[language=ATEShell]
Internal error: this window may have stopped updating. A Test Run in progress keeps recording. Details: <log path>
\end{lstlisting}

There is no dialog, because a modal box would sit in front of a run and nobody has to answer it for the
run to continue, and no abort, because the worker's record is still good. If the log write itself fails,
the Operator still gets the line, with no path claimed.

Know its limits, because they are narrow and were found by test:

\begin{itemize}
  \item it sees a raising slot (direct or queued), a timer callback, and an exception leaving a
    \texttt{QThread.run} override;
  \item it does not see a \texttt{threading.Thread} exception --- those go to \texttt{threading.excepthook},
    and the Device readers, the only such threads, catch their own;
  \item the Test Run path does not use it: \texttt{RunEngine.run}'s guard catches a worker's exception
    first and reports it through the run's own log and Test Report;
  \item \texttt{qInstallMessageHandler} was not used, because it sees Qt's messages, not Python
    exceptions;
  \item it covers the GUI launch only. \texttt{ate \_selftest} and other headless paths have no
    application log and still print.
\end{itemize}

The consequence for whoever writes a test: a child-process test can no longer trust its exit code to see
a slot exception. \texttt{assert\_no\_swallowed\_exception(stderr)} in \texttt{tests/\_fixtures.py} reads
stderr as well, and it is used by the tests that start a Qt process: the children in
\texttt{tests/test\_windows\_ci\_fixes.py}, the two in
\texttt{tests/test\_modals\_never\_block\_headless.py}, the one in
\texttt{tests/test\_gui\_launch\_30a.py}, and the packet 10a tests.

\cando{say where an unhandled exception lands, and write a test that cannot pass by swallowing one.}

\chapter{Packaging, release CI, and how a bundle is proven}\label{ch:packaging}

Four artifacts ship: a Windows installer, a Linux \texttt{.deb}, a tarball of the frozen bundle, and a
wheel. The first three come from a PyInstaller bundle; the wheel does not contain Qt at all --- it
declares the dependency and pip fetches it. Understanding which artifact carries what is the point of
this chapter, because a bug that only shows in one of them is invisible until a Station installs it.

\section{The bundle}

\texttt{packaging/\allowbreak{}build\_\allowbreak{}pyinstaller.py} builds the frozen application. It makes a throwaway venv,
installs into it, runs \texttt{packaging/\allowbreak{}pyinstaller/\allowbreak{}ate.spec}, and can smoke-test the result:

\begin{lstlisting}[language=ATEShell]
python packaging/build_pyinstaller.py --smoke
python packaging/build_pyinstaller.py --verify dist/ate   # check an existing bundle: no build, no run
\end{lstlisting}

\texttt{verify\_bundle()} is the gate a bundle must pass: it refuses a bundle that does not carry the
licence files (\texttt{\_internal/licences}, four files generated at build time by
\texttt{packaging/\allowbreak{}licence/\allowbreak{}notices.py --out build/\allowbreak{}licences}), the report template, the bundled Lua
modules, the WebEngine binary and its \texttt{.pak} and \texttt{icudtl.dat} resources --- the gate
checks that some are there, not how many. The excludes list in \texttt{ate.spec}
has its own history: PySide6's \texttt{QtWebEngineWidgets} imports \texttt{PySide6.QtWebChannel} on its
own, and keeping it excluded produced a bundle that looked fine and then died in
\texttt{ate \_pdf} with ``could not import module 'PySide6.QtWebChannel'''. It is out of the list, and
\texttt{test\_\allowbreak{}the\_\allowbreak{}spec\_\allowbreak{}does\_\allowbreak{}not\_\allowbreak{}exclude\_\allowbreak{}what\_\allowbreak{}webengine\_\allowbreak{}imports\_\allowbreak{}itself} runs the WebEngine import and
compares what it pulls in with the excludes list --- so the list cannot drift again without failing.

The real build is gated: \texttt{ATE\_\allowbreak{}BUILD\_\allowbreak{}TESTS=1 pytest -m build} runs
\texttt{tests/\allowbreak{}test\_\allowbreak{}pyinstaller\_\allowbreak{}build\_\allowbreak27a.py}, which takes minutes. It is not run by the default suite,
and the manual says so rather than implying it is.

\section{Linux: the \texttt{.deb}, the tarball, the wheel}

\texttt{packaging/\allowbreak{}linux/\allowbreak{}build\_\allowbreak{}linux.py --build --dist dist --out dist} produces the \texttt{.deb} and the
tarball from a built bundle. The \texttt{.deb} declares \texttt{DEPENDS}, the list of system libraries the
bundle cannot carry itself; \texttt{WHEEL\_SYSTEM\_LIBS} is the list for the wheel, which ships no Qt and
therefore needs the system to provide the xcb and GL libraries QtWebEngine links against. Six xcb
libraries (cursor, icccm, keysyms, shape, xkb, xkbcommon-x11) were added to that list after a bare
\texttt{debian:stable} container aborted with ``Could not load the Qt platform plugin "xcb"''.

\texttt{packaging/\allowbreak{}linux/\allowbreak{}install.sh --prefix /\allowbreak{}opt/\allowbreak{}ate --bin /\allowbreak{}usr/\allowbreak{}bin/\allowbreak{}ate} is the tarball's installer;
\texttt{packaging/\allowbreak{}linux/\allowbreak{}deb/\allowbreak{}postinst} and \texttt{postrm} are the \texttt{.deb}'s hooks (the \texttt{.deb}
is idempotent about the \texttt{ate} group; \texttt{apt purge} removes the config file).

The wheel is built with \texttt{build}. \texttt{--no-isolation} uses the interpreter already present
instead of making a throwaway venv to download setuptools into, which is what a machine with no
network to PyPI needs:

\begin{lstlisting}[language=ATEShell]
python -m build --wheel --no-isolation -o dist
\end{lstlisting}

Its version comes from the git tag, computed once by \texttt{src/ate/\_version.py} so build time and run
time agree; untagged builds fall back to \texttt{0.0.0-dev} and a build never fails because git is
unavailable. Package data matters: a bare-key bug in \texttt{pyproject.toml} meant \texttt{no wheel had
built since packet 06b}, and \texttt{tests/\allowbreak{}test\_\allowbreak{}wheel\_\allowbreak{}build\_\allowbreak28a.py} now reads the parsed config rather
than the source text, so a non-\texttt{.py} file the wheel must ship cannot go missing silently.

\section{Windows}

\texttt{packaging/\allowbreak{}windows/\allowbreak{}ate.iss} is the Inno Setup script;
\texttt{packaging/\allowbreak{}windows/\allowbreak{}build\_\allowbreak{}installer.py} compiles it. On Linux the only thing you can do is see the
command:

\begin{lstlisting}[language=ATEShell]
python packaging/windows/build_installer.py --dry-run --dist dist
\end{lstlisting}

which prints the exact \texttt{iscc} line the CI will run. \texttt{tests/\allowbreak{}test\_\allowbreak{}inno\_\allowbreak{}installer\_\allowbreak27b.py}
parses \texttt{ate.iss} and asserts against that command, and
\texttt{tests/\allowbreak{}test\_\allowbreak{}release\_\allowbreak{}ci\_\allowbreak28c.py} checks that the flags the CI passes to the installer
(\texttt{packaging/\allowbreak{}windows/\allowbreak{}smoke\_\allowbreak{}installer.py}'s \texttt{INSTALL\_FLAGS}) are the ones the manuals tell an
administrator to use.

\section{The manuals, and where they end up}

Three manuals are built before anything is packaged: the test writer's, the station administrator's
and this one. \texttt{docs/\allowbreak{}manuals/\allowbreak{}build.py} compiles each with \texttt{latexmk} in a TeX Live
image pinned by digest --- \texttt{latest} moves, and a release has to be rebuildable --- stamps the
version being released on the title page (the string \texttt{ate --version} prints, from
\texttt{src/\allowbreak{}ate/\allowbreak{}\_\allowbreak{}version.py}), and writes the PDFs into
\texttt{src/\allowbreak{}ate/\allowbreak{}manuals}, which \texttt{.gitignore} keeps out of the repository. That one staged
directory is what makes them ship: \texttt{ate.spec} sweeps every non-\texttt{.py} file under
\texttt{src/ate}, and \texttt{[tool.setuptools.package-data]} lists \texttt{*.pdf}, so the frozen bundle and
the wheel carry the same files. \texttt{verify\_bundle()} then refuses a bundle whose PDF is missing or
empty when \texttt{ATE\_\allowbreak{}RELEASE\_\allowbreak{}BUILD=1} is in the environment, which is what the release jobs set;
a build from a plain checkout says plainly that it has no manuals and carries on.

The names come from one list, \texttt{MANUAL\_\allowbreak{}NAMES} in
\texttt{docs/\allowbreak{}manuals/\allowbreak{}build.py}, and \texttt{packaging/\allowbreak{}manuals.py} reads it. A directory with a
\texttt{main.tex} that is not in the list makes the build fail, so a fourth manual cannot be left out
of the artifacts silently. The \texttt{.deb} also copies the PDFs to
\texttt{/usr/share/\allowbreak{}doc/ate}, where Debian expects documentation, and \texttt{ate --help} prints the
path the running install really has: \texttt{\_internal/ate/manuals} inside a frozen bundle,
\texttt{ate/manuals} inside the wheel's package, \texttt{/usr/share/doc/ate} in a \texttt{.deb}. On a tag
the \texttt{release} job attaches all six files --- the three PDFs and the three Markdown manuals ---
to the GitHub Release, each asset named with the version it belongs to, taken from the artifacts this
same run built and smoke-tested.

\section{The CI that runs all of this}

\texttt{.github/\allowbreak{}workflows/\allowbreak{}ci.yml} runs the suite as a single job with a
$2 \times 2$ matrix --- \texttt{ubuntu-latest} and \texttt{windows-latest}, each on CPython 3.13 and
3.14 --- so four job runs, \texttt{fail-fast: false}, with
\texttt{QT\_\allowbreak{}QPA\_\allowbreak{}PLATFORM=offscreen} and \texttt{PYTHONFAULTHANDLER=1} set at job level. Two things in it
are load-bearing and were added because of failures, not taste:

\begin{itemize}
  \item \texttt{timeout-minutes} is 30 on Linux and 40 on Windows. The first real run of this workflow
    hung on \texttt{windows-latest} for an hour on a modal dialog waiting for an Operator who was not
    there. A bounded job that reports a timeout plus a \texttt{faulthandler} stack is a job you can read.
  \item A step checks that \texttt{PySide6-Addons}, \texttt{PySide6-Essentials} and \texttt{shiboken6} are
    the same version and that the Qt the runtime reports matches them. Two Qt builds in one process is how
    the first Windows run died with \texttt{0xc0000139} (STATUS\_ENTRYPOINT\_NOT\_FOUND) the moment
    WebEngine was imported: pip had put a Qt 6.12 core next to a Qt 6.11 WebEngine.
\end{itemize}

\texttt{.github/\allowbreak{}workflows/\allowbreak{}release.yml} runs on a tag or by hand: \texttt{manuals} builds the PDFs
first, \texttt{windows-installer} and \texttt{linux-build} take them as input, and \texttt{linux-smoke}
(matrix: \texttt{deb}, \texttt{tar}, \texttt{wheel}) installs each artifact in a fresh \texttt{debian:stable}
container --- four job definitions, six job runs. On a tag only, a fifth job, \texttt{release}, creates
the GitHub Release for that tag out of the artifacts this same run uploaded. The smoke scripts are
\texttt{packaging/\allowbreak{}linux/\allowbreak{}smoke\_\allowbreak{}linux.py} and
\texttt{packaging/\allowbreak{}windows/\allowbreak{}smoke\_\allowbreak{}installer.py};
\texttt{packaging/\allowbreak{}linux/\allowbreak{}docker\_\allowbreak{}smoke.py} is the driver a human runs.
Each smoke run proves the manuals came with the artifact it installed: \texttt{ate --help} names the
directory, every name on the list is in it, each file starts with \texttt{\%PDF}, and each title carries
the tag's version. The \texttt{.deb} run looks in \texttt{/usr/share/\allowbreak{}doc/ate} too, and the
Windows run checks the bundle after the install and again after the upgrade.
The driver tags a fresh copy of the checkout to build against: \texttt{ATE\_\allowbreak{}SMOKE\_\allowbreak{}TAG}
names the tag, and with nothing named it is this checkout's own \texttt{git describe --tags}. The
staged manuals are then checked against that version rather than supplying it; read from the PDFs,
the expectation would be the thing being checked, and manuals left over from the last release would
make every check in the run agree with itself. A container built as \texttt{1.2.3} while the staged
manuals say something else is a version mismatch, not a packaging failure, and the run refuses it
rather than pass it.

\section{What has actually been run here}

Both gates were run for this manual, on the machine that wrote it, with \texttt{TMPDIR} pointed at a
directory on disk rather than \texttt{/tmp}. \texttt{ATE\_\allowbreak{}BUILD\_\allowbreak{}TESTS=1 pytest -m build} takes a few minutes
because it builds a frozen bundle; \texttt{ATE\_\allowbreak{}DOCKER\allowbreak{}TESTS=1 pytest -m build} takes several, because it builds the tarball
and the \texttt{.deb} and installs each in a fresh \texttt{debian:stable} container. Both pass. The first
attempt at the container gate in this environment died with \texttt{OSError: [Errno 28] No space left on
device} on a small \texttt{/tmp}, which is why the manual says where a gate ran and not merely that it
can. The timings quoted in chapter~\ref{ch:intro} are three separate runs of three
separate commands on one machine, not a comparison.

Packet 36's gates ran here the same way, and the numbers are this machine's. The three PDFs staged
into \texttt{src/\allowbreak{}ate/\allowbreak{}manuals} are just under 0.8 MB. The bundle built with them is 447 files and
489.6 MB; the same bundle with the three files taken out is 444 files and 488.8 MB, so the manuals
cost 0.17\% of a bundle. Packaged from the same bundle: the tarball gains 0.80 MB (199.8 MB to
200.6 MB), the \texttt{.deb} 1.6 MB (148.1 MB to 149.7 MB) --- the \texttt{.deb} pays twice, because the
PDFs ride in the bundle and are copied to \texttt{/usr/share/\allowbreak{}doc/ate} --- and the wheel gains 0.80 MB
(304 KB to 1.11 MB), which is 264\% of it: the wheel ships no Qt, so the manuals are the largest thing
in it.

What has never been run on
this machine is a Windows build, a Windows install, or an installer: the release runner did compile
\texttt{ate.iss} and install, run, upgrade and uninstall the installer at \texttt{v0.1.0-rc2} (run
41198438867), so those paths are proven somewhere --- just not here.

\cando{say which artifact carries which library, and name the test that proves a bundle is complete.}

\chapter{The test suite: seams, fixtures, gates, and how a timing test must be written}\label{ch:tests}

The suite is 1591 tests on Linux with the offscreen platform, plus the \texttt{examples/} harness that
CI also runs. As this manual was written it ran green in 336.28 s: 1586 passed, 5 skipped (the two
build-gate tests, the two container-gate tests, and the opt-in review-screenshot test). It is long
(about five and a half minutes on this machine) and it is the only place most of the rules in this
manual are actually proven, so a change to it is as consequential as a change to the code.

\section{Fixtures that every test inherits}

\texttt{conftest.py} is not boilerplate; each fixture in it exists because of a failure.

\begin{itemize}
  \item \textbf{One \texttt{QApplication} for the session.} This Qt 6.11 build segfaults at interpreter
    shutdown when a live \texttt{QApplication} is finalised alongside any \texttt{QCoreApplication}.
    Sharing one application means only a lone application is ever finalised. The binding's
    \texttt{QApplication} does not accept a \texttt{platform=} keyword, so the offscreen platform comes
    from \texttt{QT\_QPA\_PLATFORM}.
  \item \textbf{Widget trees and workers are torn down per test.} A \texttt{StagedRunScreen} keeps a live
    worker and a one-second timer for its whole lifetime; if it survives into the next test, its queued
    signals fire into that test's widgets and corrupt their visibility --- the \texttt{run\_view} flake.
    \texttt{\_stop\_workers} and \texttt{\_\allowbreak{}destroy\_\allowbreak{}widget\_\allowbreak{}tree} run before every test, and the widget
    tree is destroyed while Python is still alive.
  \item \textbf{Every started engine is stopped after every test.} \texttt{\_\allowbreak{}stop\_\allowbreak{}running\_\allowbreak{}engines}
    wraps \texttt{RunEngine.start} weakly, then aborts and joins anything still running. If it cannot, the
    test fails with ``a RunEngine worker outlived its test'' and the description of the stuck run
    (section below).
  \item \textbf{Files are closed, strictly.} On Linux an open file can be deleted; on Windows it cannot.
    The \texttt{\_release\_files} fixtures close every \texttt{RunStore}, \texttt{HeadlessRunSession} and
    \texttt{AppLog} a test opened, then assert that nothing under pytest's temp directory is still open.
    A handle that only Windows would complain about is caught here, on Linux.
  \item \textbf{The data directory is redirected.} Config defaults are written as
    \texttt{<data>/\dots}; the autouse fixture points \texttt{ate.config.\_\allowbreak{}platform\_\allowbreak{}data\_\allowbreak{}dir} at
    pytest's basetemp so no test writes to \texttt{/var/lib} or \texttt{C:\textbackslash ProgramData}.
  \item \textbf{Privilege is pinned.} \texttt{ATE\_ELEVATED=0} for the session, session-scoped on purpose:
    a module-scoped fixture that starts a run runs before any per-test fixture, so a function-scoped pin
    would arrive too late. A test that wants to be privileged sets \texttt{ATE\_ELEVATED=1}.
  \item \textbf{The PDF helper is a stub by default.} \texttt{\_fake\_pdf\_helper} replaces
    \texttt{ate.report.pdf\_\allowbreak{}job.pdf\_\allowbreak{}helper\_\allowbreak{}command} with a process that writes a 13-byte stub, except
    in the four tests that must prove a PDF has text. A test needing other helper behaviour patches it
    itself.
  \item \textbf{Qt's fatal messages are not swallowed.} \texttt{report\_qt\_fatal} keeps Qt's fatal
    message from dying inside pytest-qt's capture.
\end{itemize}

\texttt{tests/\_fixtures.py} holds the seams a test can use: \texttt{LoopbackPeer} and
\texttt{UdpLoopbackPeer} for Device tests (a TCP/UDP peer that records what it received),
\texttt{scratch\_directory} for temporary directories that conftest removes,
\texttt{executable\_stub} for a stand-in binary, \texttt{fake\_pdf\_helper} and
\texttt{recording\_\allowbreak{}pdf\_\allowbreak{}helper}, \texttt{wait\_\allowbreak{}for\_\allowbreak{}disconnect}, \texttt{watch\_store\_writes} (which
records every store call and its duration), \texttt{describe\_\allowbreak{}stuck\_\allowbreak{}runs},
\texttt{wait\_\allowbreak{}for\_\allowbreak{}aborted\_\allowbreak{}run}, \texttt{SlowAbortWarning} and
\texttt{assert\_\allowbreak{}no\_\allowbreak{}swallowed\_\allowbreak{}exception}.

\section{Two gates}

\begin{lstlisting}[language=ATEShell]
ATE_BUILD_TESTS=1 pytest -m build -q      # the real PyInstaller bundle, and the .deb built from it
ATE_DOCKER_TESTS=1 pytest -m build -q     # the debian:stable container smokes
\end{lstlisting}

Both are off by default, and both are marked \texttt{build}. The marker description in
\texttt{pyproject.toml} is the documented contract: a test marked \texttt{build} is skipped unless
\texttt{ATE\_BUILD\_TESTS=1}; the container tests skip unless \texttt{ATE\_DOCKER\_TESTS=1} and Docker is
present. \texttt{tests/\allowbreak{}test\_\allowbreak{}release\_\allowbreak{}ci\_\allowbreak28c.py} asserts that the gate is written the way it says, so a
container test cannot quietly start running on every CI job.

The five skips in a default run are therefore
\texttt{test\_\allowbreak{}pyinstaller\_\allowbreak{}build\_\allowbreak27a.py}'s frozen bundle,
\texttt{test\_\allowbreak{}linux\_\allowbreak{}packages\_\allowbreak28b.py}'s \texttt{.deb} built from it,
\texttt{test\_\allowbreak{}manuals.py}'s container permissions smoke,
\texttt{test\_\allowbreak{}release\_\allowbreak{}ci\_\allowbreak28c.py}'s three containers, and
\texttt{test\_\allowbreak{}finished\_\allowbreak{}run\_\allowbreak{}layout\_\allowbreak30b.py}'s review screenshots, which need
\texttt{ATE\_\allowbreak{}30B\_\allowbreak{}SHOTS}. Both gates were run for this manual; the numbers are in
chapter~\ref{ch:packaging}.

\section{Timing tests, and the rule that came out of 31e}

Packet 31e was two tests that failed intermittently on Windows CI and passed on Linux. Reading the notes
and the code together, the rule that emerged is worth applying to every new timing test:

\begin{enumerate}
  \item \textbf{Start the clock on an event, not on a sleep.} The old abort test started timing after
    \texttt{time.sleep(0.1)} with no proof the action had begun, so any stall of the worker around that
    point was counted as abort latency. \texttt{\_drive\_abort} now starts on the worker being inside
    \texttt{ATE.sleep}'s wait, or on \texttt{input\_asked} for a blocked input, and stops at
    \texttt{FINISHED}, before any read.
  \item \textbf{Separate the requirement from the expectation.} Story 83 says an abort should stop within
    0.5 s and must stop within 10 s. A test fails past 10 s and \textbf{warns} past 0.5 s, so a slow
    runner stays visible in CI's summary instead of being a false failure.
    \texttt{wait\_\allowbreak{}for\_\allowbreak{}aborted\_\allowbreak{}run} implements both, and prints the worker's timeline.
  \item \textbf{A stamp means what the code did, not what the test expected.} \texttt{record\_log} takes
    \texttt{time.monotonic()} when the line is recorded; the drain maps it to wall-clock later. A test
    that compares \texttt{at} with a deadline measures the scheduler, not the product. The fixed tests
    bracket each \texttt{record\_log} call on the worker and require \texttt{at} to fall inside the
    bracket.
  \item \textbf{Show that the bound can fail.} Every bound in these tests has a test that plants a stall
    and proves the test catches it: a 1 s injected stall gives ``abort took 1.011 s''; a busy loop of
    \texttt{1e10} iterations fails the 10 s requirement; a stamp taken at drain time fails the bracket.
    A bound nobody has seen violated is a bound nobody has tested.
  \item \textbf{Do not widen a timeout.} Packet 32b's checklist says it plainly: the 10 s wait stays and
    can still fail, and when it does it now says what the worker was doing.
  \item \textbf{Run under load.} The convention the packets use is ten repetitions of the changed files
    with six busy loops on three cores. Several flakes only appear under contention, and two shared-helper
    races (the \texttt{LoopbackPeer} that stopped reading at \texttt{stop()} and the VISA tests that read
    \texttt{peer.received} before the peer had seen the connection end) were found that way.
\end{enumerate}

\section{What \texttt{describe\_\allowbreak{}stuck\_\allowbreak{}runs} gives you}

When a run does not end, the failure message includes a report section: every live worker and session,
what is recorded in the store, the store's write timings with any write still in flight, and all thread
stacks from \texttt{faulthandler}. \texttt{pytest\_\allowbreak{}runtest\_\allowbreak{}makereport} adds it to any test that failed
with a worker running, whichever wait helper it used. Its purpose is to answer the question packet 31e
could not: whether a stall is the disk (long or in-flight writes) or something else (normal writes, with
the worker's stack elsewhere).

\section{Where a new test goes}

Name it after the packet that added it, in the style the suite already uses
(\texttt{tests/\allowbreak{}test\_\allowbreak{}device\_\allowbreak{}abort\_\allowbreak32b.py}). Put shared helpers in \texttt{tests/\_fixtures.py}, not in a
test file, so the next test does not copy a helper that has a known race. If your test starts a Qt process,
use \texttt{assert\_\allowbreak{}no\_\allowbreak{}swallowed\_\allowbreak{}exception}: exit 0 no longer means nothing went wrong.

\cando{write a timing test that can fail, and pick the right gate for a slow one.}

\chapter{PySide6 traps, found the hard way}\label{ch:pyside}

ATE uses PySide6 (ADR-0001, \texttt{docs/\allowbreak{}adr/\allowbreak0001-pyside6-instead-of-pyqt6.md}: PyQt6's wheels are GPL,
and shipping a GPL binding under an MIT label is a licence gap). The switch from PyQt6 found real
behaviour differences. Each is a trap for someone writing code or a test here.

\begin{itemize}
  \item \textbf{An exception in a slot or \texttt{QThread.run} does not end the process.} PyQt6 aborted
    with SIGABRT (exit 134). PySide6 prints the traceback, calls \texttt{sys.excepthook} and carries on.
    That is why the guard in \texttt{RunEngine.run} and the hook in \texttt{ate.exception\_log} exist ---
    not as style, but because otherwise the run row stays open and the screen waits forever.
  \item \textbf{A child process's exit code is not a signal.} A test that started a Qt process used to
    prove nothing had gone wrong by \texttt{returncode == 0}. It no longer can. Use
    \texttt{assert\_\allowbreak{}no\_\allowbreak{}swallowed\_\allowbreak{}exception(stderr)}.
  \item \textbf{\texttt{Signal(int)} is a C \texttt{int}.} \texttt{emit(2**31)} with a slot connected
    returns normally and the slot receives -2147483648 with a \texttt{RuntimeWarning}; the
    \texttt{OverflowError} is raised only when nothing is connected, and it can surface later as a
    \texttt{SystemError} in a slot's next C call. \texttt{Signal(str)} turns \texttt{5} and \texttt{None}
    into \texttt{''}. A slot cannot see a bad value, so validate on the sending side. Declare
    \texttt{Signal(object)} for anything that can exceed \texttt{2**31 - 1}.
  \item \textbf{Two Qt builds in one process.} PySide6-Addons pins PySide6-Essentials and shiboken6 to the
    exact same version, and CI checks that the Qt the runtime reports matches the bindings. If it does not,
    WebEngine loads a second Qt and Windows dies with \texttt{0xc0000139} the moment WebEngine is
    imported. On Linux it tolerates it --- until it does not.
  \item \textbf{\texttt{QtWebEngineWidgets} imports \texttt{QtWebChannel}.} Excluding it from the bundle is
    the packaging bug that produced a bundle that looked complete and failed at render time.
  \item \textbf{Deleted objects: different message, same handling.} PyQt6 raises
    \texttt{wrapped C/\allowbreak{}C++ object of type Thing has been deleted}; PySide6 raises
    \texttt{libshiboken: Internal C++ object (Thing) already deleted.} Every catch here is
    \texttt{except RuntimeError}, so both behave the same. \texttt{sip.isdeleted(x)} became
    \texttt{not shiboken6.isValid(x)}.
  \item \textbf{A second \texttt{QCoreApplication} raises.} Creating one while a \texttt{QApplication}
    exists raises \texttt{RuntimeError} in PySide6 and was tolerated by PyQt6. This is why the suite keeps
    exactly one application --- and why \texttt{conftest.py} destroys the widget tree before quitting: an
    orphaned \texttt{QWidget} finalised at interpreter shutdown segfaults in this Qt 6.11 build either way.
  \item \textbf{Small API differences that only show in tests.} \texttt{findChildren} takes one type, not a
    tuple. \texttt{QWidget.cursor()} does not compare equal to a \texttt{Qt.CursorShape} enum: read
    \texttt{.shape()}. \texttt{QWIDGETSIZE\_MAX} is not exported, so the tests use the literal
    \texttt{(1 << 24) - 1}.
  \item \textbf{Offscreen on Windows has no fonts.} Qt's offscreen platform on Windows looks for fonts in
    Qt's own (empty) library directory instead of asking Windows, so every character draws as the
    missing-glyph box and every width is wrong. \texttt{conftest.py} sets \texttt{QT\_QPA\_FONTDIR} to
    Windows' font directory before any Qt import, and every subprocess a test starts inherits it. A real
    operator's \texttt{windows} platform has no such problem; a test on a Windows runner does.
\end{itemize}

None of these were run ``on both bindings'': only PySide6 was tested here. The ADR records the
comparison as a reviewer's probe where a claim rests on one.

\cando{name a difference between PySide6 and PyQt6 that changes how a test must be written.}

\chapter{How to read a failed Test Run}\label{ch:failed-run}

A run that went wrong leaves four traces. Read them in this order, because the order is what makes the
cause visible.

\section{The run row first}

\texttt{test\_run.result} is \texttt{PASS} or \texttt{FAIL}. \texttt{end\_reason} is the story:

\begingroup\small
\begin{longtable}{p{3.2cm} p{9.4cm}}
\toprule
\textbf{End Reason} & \textbf{What happened} \\
\midrule
\endhead
\texttt{COMPLETED} & the schedule ran to its end, abort or not \\
\texttt{PRETEST\_FAILED} & a Pre-test Action raised; no step ran, and every step is NOTRUN \\
\texttt{STOPPED\_ON\_FAIL} & the Operator answered \texttt{ABORT} at a Stop-on-Fail prompt \\
\texttt{ABORTED} & a run-level \texttt{request\_abort()}, wherever it landed --- at a boundary, inside a
  step, or at a pause (packet 32b: the same button now gives the same reason) \\
\texttt{INTERRUPTED} & the worker ended on an exception it could not handle, or the process died with the
  run open and launch closed it \\
\bottomrule
\end{longtable}
\endgroup

A step's \texttt{result} is one of \texttt{PASS}, \texttt{FAIL}, \texttt{ERROR}, \texttt{SKIPPED},
\texttt{NOTRUN}, \texttt{FORCED\_NOTRUN} or \texttt{NOREPORT}. \texttt{NOREPORT} means the step ran and never
called \texttt{ATE.report} --- a schedule bug, recorded as a WARNING in the Run log rather than silently
counted as a pass.

\section{The two logs}

The \textbf{Run log} is the \texttt{log\_line} table: one row per line, with a level
(\texttt{DEBUG}/\texttt{INFO}/\texttt{WARNING}/\texttt{ERROR}), a source (\texttt{schedule}, \texttt{ate} or
\texttt{device}) and an owner --- the Action or the Attempt that produced it. Run log lines come from
\texttt{ATE.log}, redirected \texttt{print}, Device failures and warnings, and trace lines when tracing is
on. The \textbf{application log} is the per-launch file at \texttt{<data>/\allowbreak{}logs/\allowbreak{}ate-app-<time>.log}, which
is where an unhandled exception's traceback goes. They are different things and answer different questions:
``what did this run do'' versus ``what did the application do wrong''.

The last line of a step's Attempt in the report is the last \texttt{ATE.report} call of that Attempt; a
step whose Attempt raised records the error text as its \texttt{msg}.

\section{When the failure is the test harness, not the run}

If a test fails with ``a RunEngine worker outlived its test'' or ``N RunEngine worker(s) did not stop'',
the report section added by \texttt{describe\_\allowbreak{}stuck\_\allowbreak{}runs} is the thing to read: the live workers, what each
session thinks its state is, what the store has recorded, the store's write timings with any write still in
flight, and every thread's stack. Long or in-flight writes mean the disk. Normal writes with the worker's
stack elsewhere mean a fault the suite has not found yet. That is the diagnostic packet 32b added because
the Windows stall it could not reproduce had no explanation.

Read the warnings that follow an abort carefully: ``left open at end of run'' means a Device the schedule
never closed; ``was still closing when the aborted run ended'' means the close outlived the run's 3 s bound
and will finish later; ``datagram truncated to N byte(s)'' and ``timed out after N byte(s)'' are different
read failures, and the Device says which one it is.

\section{Reading the record with a query}

The fastest way to see what a run recorded, from a Station's database:

\begin{lstlisting}[language=ATEShell]
python -c "import sqlite3, sys; rows=sqlite3.connect(sys.argv[1]).execute('SELECT step_number,result,attempt_count FROM node ORDER BY node_seq'); print(rows.fetchall())" /tmp/station/data/ate.db
\end{lstlisting}

Read-only, no ATE involved. The database is passed as an argument rather than written inside the
string: on Windows the path has backslashes, and a path pasted into a string literal in
\texttt{-c} code is read as an escape sequence --- \texttt{C:\textbackslash Users\textbackslash ...}
is a \texttt{SyntaxError} before the query ever runs.

The same query against \texttt{log\_line} ordered by \texttt{seq} is the Run log
in emission order, which is the order the report shows.

\cando{go from a failed run to the row, the log line and the thread that produced it.}

\chapter{Worked example: adding a Device Type}\label{ch:device-type}

This chapter is not a description of a Device Type. It is a change made to this codebase, and every
line in it was executed: \texttt{tests/test\_manuals.py} copies \texttt{src/} into a temporary
directory, puts the driver below at \texttt{src/ate/devices/echo.py}, applies the three edits shown
here character for character, and runs \texttt{docs/manuals/examples/developer/echo.tslua} headless
against that copy. If the manual and the code ever part company, that test fails.

A Device Type has three parts, and all three have to exist: the driver that opens it, the registry
line that makes the name mean that driver, and the enum member that lets a schedule write the name
down. Two of them are in the same file.

\section{Step 1: the driver}

A driver is a subclass of :class:`ate.devices.registry.Driver`. It declares the Device Type name,
the fields an \texttt{ATE.open} may pass, and the fields \texttt{ATE.read}/\texttt{ATE.write} may
pass; the registry validates every one of them generically through :func:`validate_params`, which
raises a :class:`ValueError` naming the offending field before anything connects.

\lstinputlisting[language=ATEPython, caption={\texttt{docs/manuals/examples/developer/echo_driver.py}}]{../examples/developer/echo_driver.py}

Three things in that file are not decoration:

\begin{itemize}
  \item \texttt{fields} is the whole of the open-time contract. An undeclared key is refused with
    \texttt{unknown field 'x' for ECHO device}, so a schedule typo never reaches a driver.
  \item \texttt{read\_fields} and \texttt{write\_fields} are \texttt{READ\_FIELDS} and
    \texttt{WRITE\_FIELDS} from \texttt{src/ate/devices/stream.py}. Every Device Type takes the same
    read and write fields; a Device with a different contract would declare its own dict, and the
    ATE wrappers validate against the driver carried on the handle.
  \item \texttt{partial\_read\_is\_truncation = True} on the handle is what makes a cut-short read
    say \texttt{datagram truncated to N byte(s)} instead of \texttt{timed out after N byte(s)}. The
    engine asks the Device which it is (packet 32b review): there is no stream and no deadline here,
    so claiming a timeout would be a lie.
\end{itemize}

\section{Step 2: register the driver}

The registry is built once per run, in \texttt{LuaHost.\_\_init\_\_}
(\texttt{src/ate/engine/lua_host.py}). Add one line beside the others:

\begin{lstlisting}[language=ATEPython]
        self._devices.register("TCP", TcpDriver)
        self._devices.register("SERIAL", SerialDriver)  # packet 18a: pyserial ports, same API
        self._devices.register("UDP", UdpDriver)  # packet 18b: datagrams, one datagram per message
        self._devices.register("VISA", VisaDriver)  # packet 19: pyvisa resources, read on demand (no reader)
        self._devices.register("ECHO", EchoDriver)  # the developer manual's worked example
\end{lstlisting}

and the import above it:

\begin{lstlisting}[language=ATEPython]
from ..devices import DeviceRegistry, EchoDriver, TcpDriver
\end{lstlisting}

Without this line the enum member is a name with nothing behind it: \texttt{ATE.open(DeviceType.ECHO,
\{label = "x"\})} reaches the registry, which raises \texttt{unknown DeviceType 'ECHO'}. That test is
run too --- see the last section of this chapter.

\section{Step 3: the enum member}

The flat enum tables a schedule can see live in \texttt{\_ENUMS} in the same file:

\begin{lstlisting}[language=ATEPython]
    "DeviceType": {
        "TCP": {"TCP": "TCP"},
        "SERIAL": {"SERIAL": "SERIAL"},
        "UDP": {"UDP": "UDP"},
        "VISA": {"VISA": "VISA"},
        "ECHO": {"ECHO": "ECHO"},
    },
\end{lstlisting}

Each entry maps the Lua table name to its members; :func:`_device_type_name` (also in
\texttt{lua_host.py}) turns a member back into the string the registry looks up, and it is the one
that raises \texttt{deviceType must be a DeviceType enum member} when a schedule passes a bare
string or a number.

\section{Step 4: run it}

The schedule below is the one the test runs. It writes a terminator, leaves a Device open at the end
of the run, and cuts a read with \texttt{size} --- three things a driver has to get right, and each
of them shows up in the record.

\lstinputlisting[language=ATELua, caption={\texttt{docs/manuals/examples/developer/echo.tslua}}]{../examples/developer/echo.tslua}

To run it yourself against a copy of the source with the driver in place, put the runner from
chapter~\ref{ch:run-session} in \texttt{runner.py} and start it with the copy of \texttt{src} on
\texttt{PYTHONPATH}:

\begin{lstlisting}[language=ATEShell]
python runner.py /tmp/station/config/config.json "Developer: echo Device"   # PYTHONPATH=<the copy of src>
\end{lstlisting}

\section{What the run proves}

The test asserts the run ends \texttt{PASS} / \texttt{COMPLETED} with all three steps
\texttt{PASS}, and then reads the traces the engine recorded:

\begin{center}
\small
\begin{tabular}{p{5.4cm} p{7.6cm}}
\toprule
\textbf{Line in the Run log} & \textbf{What it proves} \\
\midrule
\texttt{open ECHO manual-example ok (0.0 ms)} & the registry wrapped the new driver like any other,
  and traced it \\
\texttt{write ECHO manual-example 6 bytes "hello\textbackslash n" (0.0 ms)} & the terminator is part
  of the counted payload \\
\texttt{read ECHO sized 4 bytes "ABCD" partial (0.0 ms)} & \texttt{size} cut the read short, and the
  handle reported it as partial \\
\texttt{ATE.read: datagram truncated to 4 byte(s)} & the Device said it was a truncation, not a
  timeout \\
\texttt{device ECHO left-open left open at end of run} (WARNING) & the run-end teardown closed a
  Device the schedule never closed \\
\bottomrule
\end{tabular}
\end{center}

Two more assertions are the negative cases: a Device Type the registry has never registered raises
\texttt{unknown DeviceType 'ECHOX'} from \texttt{DeviceRegistry.open\_device}, and the schedule run
with only the enum member added and no register line fails the run with \texttt{deviceType must be a
DeviceType enum member}. A Device Type is only real when all three parts are there.

\section{What else a new Device Type touches}

\begin{itemize}
  \item A test file named after the packet that added it
    (\texttt{tests/test\_device\_abort\_32b.py} is the style), with the
    \texttt{LoopbackPeer}/\texttt{UdpLoopbackPeer} fixtures in \texttt{tests/\_fixtures.py} if the
    Device needs a peer.
  \item The Device Type list in both test-writer manuals:
    \texttt{missing\_from(MANUALS["test-writer"].texts(), device\_types())} is derived from the
    drivers that exist, so a driver nobody documented is a failing suite, not an omission somebody
    might notice later.
  \item \texttt{docs/manuals/llm/test-writer.md}: a schedule writer reading only that file needs to
    know the Device exists and what its fields accept.
\end{itemize}

\cando{add a Device Type with a driver, a register line and an enum member, and show it in the record.}

\chapter{Worked example: adding a schema version}\label{ch:schema-version}

The Test Run record is defined in one file, \texttt{src/\allowbreak{}ate/\allowbreak{}storage/\allowbreak{}schema.py}, and one number:
\texttt{SCHEMA\_VERSION}. Everything that reads or writes the record goes through
\texttt{RunStore}, so a table added here is the only way a table exists. The example below is run by
\texttt{tests/\allowbreak{}test\_\allowbreak{}manuals.py}: it migrates a database with the shipped code,
tries a copy of the code that bumps only the number, migrates the same database with a copy that
declares version 6 and one more table, and then shows the shipped code refusing that database.

\section{The two edits}

Add the table to the metadata, and bump the version. The comment that explains each bump is part of
the convention: the file's own history is written in those comments --- \texttt{input\_value} made it
2, \texttt{media\_release} and \texttt{run\_file} made it 3, \texttt{schedule\_instance} made it 4,
\texttt{equipment} made it 5.

\begin{lstlisting}[language=ATEPython]
SCHEMA_VERSION = 6
\end{lstlisting}

\begin{lstlisting}[language=ATEPython]
review_note = Table(
    "review_note",
    metadata,
    Column("run_uuid", Text, primary_key=True),
    Column("seq", Integer, primary_key=True),  # the composite key, as every other run table
    Column("text", Text, nullable=True),
    Column("at", Text, nullable=True),
)
\end{lstlisting}

Nothing else is required: no hand-written DDL, no migration script. The version lives in
\texttt{PRAGMA user\_version}, and the tables are created by
\texttt{metadata.create\_\allowbreak{}all(engine)} in \texttt{RunStore.\_\allowbreak{}ensure\_\allowbreak{}schema}, which reflects each
table through the backend's own catalog and so never raises on a database that already has it.

The types are the ones the record uses everywhere: UUIDs as 36-character \texttt{Text}, times as
ISO-8601 \texttt{Text}, booleans as \texttt{Integer} 0/1, structured values as JSON \texttt{Text}.
There is deliberately no SQLite-specific column: a future Postgres backend is a connection change,
not a schema change. The only raw SQL in the storage layer is the two pragmas (WAL mode and
\texttt{user\_version}), and they stay in \texttt{store.py}.

\section{What a migration does, and what it does not}

\begin{lstlisting}[language=ATEShell]
python -m ate db migrate /tmp/station/data/ate.db --yes
\end{lstlisting}

prints \texttt{backed up \dots to \dots.backup-<YYYYMMDD-HHMMSS>} and then
\texttt{created or upgraded database at \dots (schema version 6)}. In the test this is run against the
same file more than once: with the shipped code it creates the ten tables the record has and makes no
copy, because there was no database to copy; with the patched copy it takes the version 5 database as
a copy beside itself, then adds \texttt{review\_note} on top. The assertion that matters is that the
set of tables only grew --- a migration adds, it never rewrites or drops --- and that the copy is the
database as it was, version 5 and no \texttt{review\_note}.

The last run is the reason the version check exists: the shipped code, pointed at the version 6
database, exits 1 with \texttt{this ATE is too old for \dots; update ATE}. \texttt{check\_\allowbreak{}database\_\allowbreak{}at\_\allowbreak{}launch}
compares \texttt{PRAGMA user\_version} with \texttt{SCHEMA\_VERSION} and returns
\texttt{NEEDS\_MIGRATE} when the database is behind, \texttt{TOO\_OLD} when it is ahead, and
\texttt{OK} only when they match. A run against a database whose tables are missing fails at the
first insert; the check is what makes the failure a message instead of an \texttt{OperationalError}
in a worker.

Spec rule 20 asked for a timestamped backup before an upgrade. It is no longer something the
Administrator has to do by hand: \texttt{ate db migrate} and \texttt{ate setup} both go through
\texttt{src/ate/storage/backup.py}, which copies the database beside itself before a migration runs.
The admin manual's database chapter says what an Administrator still has to back up themselves.

\section{Adding a column instead of a table}

\texttt{create\_\allowbreak{}all} does not alter a table that already exists, so adding a column to
\texttt{node} with a version bump leaves existing databases without it. Packet 35 closed the quiet
version of that mistake, the one where nothing at all is created: \texttt{RunStore.\_ensure\_\allowbreak{}schema}
stamps \texttt{PRAGMA user\_\allowbreak{}version} only when \texttt{missing\_\allowbreak{}tables} found something to
create. A database that stores an older version but already holds every table this ATE declares gets
no stamp and no silent pass. It has its own launch status, \texttt{NO\_\allowbreak{}MIGRATION}, and
\texttt{ate db status}, \texttt{ate db migrate} and the launch dialog refuse it with the same line:

\begin{verbatim}
the database at /var/lib/ate/ate.db stores schema version 5 and already holds
every table this ATE declares: this ATE has no migration from version 5 to 6,
so it is not stamped as version 6, and `ate db migrate` on it changes nothing
\end{verbatim}

The test runs exactly that: the shipped code, with only the number changed, exits 1 with that
sentence, and the database still says version 5 afterwards. Nothing in
\texttt{src/ate} writes an \texttt{ALTER} or an \texttt{add\_column}, and every insert names its columns
explicitly (\texttt{insert\_test\_run} builds its row field by field), so a new column has to be written
and tested as an explicit migration. The version number is a guard, not a migration.

What packet 35 did not close is the case the guard cannot see. A version bump that adds a column
\emph{alongside} a new table is still stamped silently: \texttt{missing\_tables} finds the new table, so
the stamp is raised, and the column on the table that already existed is never added. The database
then claims version 6 with an old \texttt{node} in it, and nothing notices. Before the next schema
version there has to be a migration registry --- a list of steps, each from version N to N+1, run in
order, each tested against a database at N --- so that a version number means "these steps ran"
instead of "some table was created".

\section{Where a new field reaches the report}

A table is not a report. The HTML writer in \texttt{src/\allowbreak{}ate/\allowbreak{}report/\allowbreak{}html.py} reads the record through
the store's query methods; a new table that nothing queries is invisible on the printed page. The
batching rules in chapter~\ref{ch:storage} apply to a new row type as well:
\texttt{\_\allowbreak{}rows\_\allowbreak{}per\_\allowbreak{}statement} keeps an insert under SQLite's variable limit, and
\texttt{\_log\_seq} is what orders a new row against the rest of the run.

\cando{add a table, bump the version, and show what a database at a different version does.}

\chapter{Documents, ADRs, and where a change has to go}\label{ch:docs}

There are five kinds of document in this repository and they are not equal. When they disagree, the
code is the record: a document that says something the code does not do is a bug in the document, and
the manual chapters in this packet are written from the code, not from the documents.

\begingroup\small
\begin{longtable}{p{4.4cm} p{8.2cm}}
\toprule
\textbf{Where it is} & \textbf{What it is good for} \\
\midrule
\endhead
\texttt{GLOSSARY.md} & the words. \texttt{Test Run}, \texttt{Test Step}, \texttt{Attempt},
  \texttt{Media Release}, \texttt{Device Type}: every manual and every test uses these, so a new
  term has to be added here before it is used anywhere \\
\texttt{docs/adr/} & decisions that were argued and settled. ADR-0001 is PySide6 rather than PyQt6
  (the licence question). If you change a binding, a runtime or a packaging choice, an ADR is where
  the reason goes \\
\texttt{src/ate/} and \texttt{tests/} & the implemented behaviour and its executable checks. When
  written guidance and these checks differ, follow the code and tests \\
\texttt{docs/manuals/} & the three manuals and the three Markdown manuals, held to the code by
  \texttt{tests/\allowbreak{}test\_\allowbreak{}manuals.py} \\
\texttt{docs/\allowbreak{}building-from-source.md} & how to build, including how to build these PDFs \\
\bottomrule
\end{longtable}
\endgroup

\section{Where a change has to go, by kind}

A change is not finished when the code works. Each kind of change has documents that follow it, and
the suite checks most of them:

\begin{itemize}
  \item \textbf{A new config key} --- \texttt{src/ate/config.py}'s \texttt{\_SCHEMA} and the admin
    manual, both LaTeX and Markdown. \texttt{config\_keys()} is derived from \texttt{\_SCHEMA}, so a
    key added to the code and not to the manual fails the suite.
  \item \textbf{A new \texttt{ate} subcommand or flag} --- \texttt{src/ate/cli.py} and the admin
    manual. \texttt{cli\_names()} reads the parser itself, including the hidden
    \texttt{\_pdf}/\texttt{\_selftest} helpers and the \texttt{db}/\texttt{definition} subcommands
    the parser never sees.
  \item \textbf{A new \texttt{ATE} function, enum or setter} --- \texttt{lua\_host.py} and both
    test-writer manuals. Same mechanism, same failure.
  \item \textbf{A new Device Type} --- the driver, the register line, the enum member, and both
    test-writer manuals (\texttt{device\_types()} and \texttt{device\_drivers()} derive it from the
    classes that exist). Chapter~\ref{ch:device-type} is the worked version.
  \item \textbf{A new module under \texttt{src/ate}} --- the developer manual and
    \texttt{docs/\allowbreak{}manuals/\allowbreak{}llm/\allowbreak{}developer.md}. \texttt{ate\_modules()} walks the package, so a module
    nobody mentioned is a red suite.
  \item \textbf{A new table in the record} --- \texttt{schema.py}, the version bump, and both
    developer manuals (\texttt{schema\_tables()} reads \texttt{metadata.tables}).
  \item \textbf{A new command a manual tells someone to run} --- it goes into
    \texttt{COMMANDS\_RUN} or \texttt{COMMANDS\_NOT\_RUN} (admin) or \texttt{DEV\_COMMANDS\_RUN} or
    \texttt{DEV\_\allowbreak{}COMMANDS\_\allowbreak{}NOT\_\allowbreak{}RUN} (developer), with the exit code and the output it must produce,
    or a reason it is not run. A command in a listing that nobody runs and nobody explains fails.
  \item \textbf{A change to the staged screens, the abort path or the timing of a run} --- the ADRs
    if the decision changes, \texttt{GLOSSARY.md} if a term changes, and the packet file if a rule is
    being bent.
\end{itemize}

\section{Writing an ADR}

One decision per file, named \texttt{NNNN-short-sentence.md}, with the context, the decision, the
consequences, and what was not tested. ADR-0001 is the pattern: it records why PySide6 (MIT) was
chosen over PyQt6 (GPL), and it is where the licence reason lives so that nobody re-litigates it in
a pull request. Where a claim rests on a comparison that was never run on both sides --- the
PySide6/PyQt6 exception behaviour, for instance --- the ADR says so. That is the difference between a
decision record and a marketing page.

\section{Reading implementation notes as a developer}

The source and executable tests are the public record of current behaviour. The manual chapters
call out limits that have not been verified on the target platform; the Windows-only Qt and
installer behaviour still needs a real Windows run.

\cando{say which document a given change must also update, and which document is not authoritative.}

\chapter{What this manual cannot state, and where the documents are wrong}\label{ch:unproven}

A manual that does not say what it has not tested is a manual a developer trusts in the wrong place.
This chapter is the list of things this manual asserts only as ``the code says'', not as ``we saw it
happen''.

\section{Never run here}

\begin{itemize}
  \item \textbf{Real hardware.} Every Device test in the suite talks to something in the same
    process: the \texttt{LoopbackPeer} and \texttt{UdpLoopbackPeer} in \texttt{tests/\_fixtures.py},
    a socket on an ephemeral port, a pseudo-terminal for \texttt{SERIAL}, and pyvisa-py behind a
    \texttt{TCPIP0::127.0.0.1::\dots::SOCKET} resource for \texttt{VISA}. No GPIB instrument, no real
    serial adapter, no NI-VISA stack has ever been opened by this suite. The driver contract is
    proven; the driver against an instrument is not.
  \item \textbf{Real Windows sessions.} The same tests run on \texttt{windows-latest} in CI, and
    the notes record those CI runs being green. Nothing in this manual was executed on Windows by
    us. The \texttt{icacls} permission paths, the Inno installer and the frozen Windows bundle are
    asserted by reading the scripts and by \texttt{tests/test\_release\_ci\_27c.py}, which checks the
    workflow's shape, never by an install.
  \item \textbf{The Debian container smokes, on this machine.}
    \texttt{ATE\_DOCKER\_TESTS=1 pytest -m build} and \texttt{packaging/linux/docker\_smoke.py} were
    attempted while writing this manual and did not complete: the container ran out of scratch space.
    \texttt{tests/test\_release\_ci\_28c.py} asserts the gate is written as documented and the smoke
    scripts' logic; the containers themselves are CI's job. Nothing in this manual claims those
    smokes passed here.
  \item \textbf{A real Chromium PDF.} \texttt{conftest.py} replaces the PDF helper with a stub for
    every test except the four that must prove a PDF has text. The sandbox decision, the flags and
    the failure messages are tested; the rendered page is not, except in those four.
  \item \textbf{Both Qt bindings.} The PySide6/PyQt6 differences in chapter~\ref{ch:pyside} come
    from the 32a switch and its notes. They were not measured side by side on this machine; where a
    claim rests on one binding only, ADR-0001 says so.
\end{itemize}

\section{The 32b stall: fixed, not explained}

Packet 32b is the clearest example of a fix that outran its diagnosis. One Windows CI run left a
Test Run stalled for tens of seconds. The packet added the abort-boundary handling, the
\texttt{describe\_stuck\_runs} report, and the store's write timing, and the tests prove each of
them behaves. The cause --- that the disk was the thing that was slow --- is not established. The
diagnostic exists so that the next stall can be read: long or in-flight writes point at the disk,
normal writes with the worker's stack elsewhere point at something nobody has found. A manual that
stated the cause as fact would be the same kind of error as a test that widens a timeout to make a
failure disappear.

\section{Where the documents and the code disagree}

These were found while writing the manuals. In every case the code was followed and the document was
left as it is, so that the difference is visible rather than smoothed over:

\begin{itemize}
  \item \textbf{Spec rule 20 says \texttt{ate db migrate} upgrades ``after a timestamped backup''.}
    It makes no backup. \texttt{cli.\_cmd\_db} calls \texttt{RunStore}, which creates the missing
    tables and sets \texttt{PRAGMA user\_version}, and prints
    \texttt{created or upgraded database at \dots}. The admin manual tells an administrator to back
    up first, and says so in the same breath as the command.
  \item \textbf{Spec rule 08 says \texttt{log\_line} rows are written ``about once a second''.}
    They are inserted per step, not on a timer. The review notes recorded this in packet 08 and it
    has never been corrected in the spec.
  \item \textbf{The release workflow has three jobs, not five.} The 32a notes and the review notes
    both say ``all 5 jobs''; \texttt{release.yml} defines \texttt{windows-installer},
    \texttt{linux-build} and \texttt{linux-smoke}, and \texttt{linux-smoke} runs a matrix of
    \texttt{deb}, \texttt{tar} and \texttt{wheel} --- five job runs from three job definitions. The
    count is right; the word is not.
  \item \textbf{\texttt{\_location()} finds a step's file and line by searching the schedule text}
    for \texttt{:setAction(}, not from the Lua call stack. For a multi-line Action, a helper module
    or a Sub-schedule the reported location can be wrong. The manual describes what is reported, not
    what the notes imply it should be.
  \item \textbf{The prototype layout in \texttt{run-screen-prototype.html} is not the screen.} The
  staged Run screen is a vertical stack; the review notes say so, and the developer manual describes
  the widgets the code creates.
\end{itemize}

\section{What changes the moment you read this}

Two numbers in this manual are snapshots: the suite's size and the elapsed time. Both were taken on
one machine on one day, and both move with every packet. The names, the signals, the tables and the
drivers are not snapshots --- the suite derives them from the code and fails if a manual stops naming
them. If \texttt{1571} is no longer the number, the tests in \texttt{tests/test\_manuals.py} do not
care; if a signal is gone from \texttt{RunEngine} and the manual still lists it, nothing notices, so
the direction of the risk matters: the checks catch a manual falling behind, not a manual ahead of
the code.

\cando{say what a claim in this manual rests on, and what would falsify it.}


\end{document}
